\documentclass[a4paper,11pt]{amsart}
\usepackage[margin=1in]{geometry}
\usepackage{amsmath,amssymb,amsthm}
\usepackage{aliascnt}
\usepackage{booktabs}
\usepackage{enumitem}
\usepackage[colorlinks=true,linkcolor=blue,citecolor=blue,urlcolor=blue]{hyperref}
\usepackage[capitalize]{cleveref}
\emergencystretch=3em

\newtheorem{theorem}{Theorem}[section]
\newaliascnt{lemma}{theorem}
\newtheorem{lemma}[lemma]{Lemma}
\aliascntresetthe{lemma}
\newaliascnt{proposition}{theorem}
\newtheorem{proposition}[proposition]{Proposition}
\aliascntresetthe{proposition}
\newaliascnt{corollary}{theorem}
\newtheorem{corollary}[corollary]{Corollary}
\aliascntresetthe{corollary}
\newaliascnt{problem}{theorem}
\newtheorem{problem}[problem]{Problem}
\aliascntresetthe{problem}
\newaliascnt{definition}{theorem}
\newtheorem{definition}[definition]{Definition}
\aliascntresetthe{definition}
\theoremstyle{remark}
\newaliascnt{remark}{theorem}
\newtheorem{remark}[remark]{Remark}
\aliascntresetthe{remark}
\crefname{theorem}{Theorem}{Theorems}
\Crefname{theorem}{Theorem}{Theorems}
\crefname{lemma}{Lemma}{Lemmas}
\Crefname{lemma}{Lemma}{Lemmas}
\crefname{proposition}{Proposition}{Propositions}
\Crefname{proposition}{Proposition}{Propositions}
\crefname{corollary}{Corollary}{Corollaries}
\Crefname{corollary}{Corollary}{Corollaries}
\crefname{definition}{Definition}{Definitions}
\Crefname{definition}{Definition}{Definitions}
\crefname{remark}{Remark}{Remarks}
\Crefname{remark}{Remark}{Remarks}
\crefname{problem}{Problem}{Problems}
\Crefname{problem}{Problem}{Problems}

\newcommand{\F}{\mathbb F}
\newcommand{\K}{\mathbb K}
\newcommand{\C}{\mathbb C}
\newcommand{\B}{\mathbb B}
\newcommand{\RS}{\operatorname{RS}}
\newcommand{\Q}{\operatorname{Q}}
\newcommand{\emca}{\varepsilon_{\mathrm{mca}}}
\newcommand{\eca}{\varepsilon_{\mathrm{ca}}}
\newcommand{\Htwo}{H_2}
\newcommand{\poly}{\operatorname{poly}}
\newcommand{\Fib}{\operatorname{Fib}}
\newcommand{\Span}{\operatorname{span}}
\newcommand{\wdeg}{\operatorname{wdeg}}
\newcommand{\Cen}{\operatorname{Cen}}
\newcommand{\LineRay}{\operatorname{LineRay}}
\newcommand{\dist}{\operatorname{dist}}
\newcommand{\wt}{\operatorname{wt}}

\title{Proof-Audited and Expanded Bounds and Reductions for RS-MCA}
\date{July 2026}

\begin{document}

\begin{abstract}
This expanded proof-audited note separates genuine theorem-level material from the still-missing dense-frontier upper ledger for smooth-domain Reed--Solomon MCA.  The unsafe side is reduced to exact integer prefix floors and the simple-pole list-to-MCA conversion.  On the safe side the note keeps only statements proved under their displayed hypotheses: broad safe anchors, quotient descent, split-pencil reductions, second-moment identities, unconditional but weak prefix/split-pair ceilings, the exact saturation identity for MCA codeword rays, and the theorem that the shift-pair ledger is downstream from Q once Q is stated as a max-fiber bound.  This revision also fixes endpoint and notation issues found in the audit, expands the finite-arithmetic tables, adds line-ray saturation bookkeeping, proves the logarithmic-moment equivalence for primitive asymptotic Q, and adds the conditional asymptotic closure theorem from the entropy-scale inverse input.  It also records a Fourier-flat sufficient condition for Q in large-characteristic leaves.  It does not claim the adjacent deployed safe rows: an attempted Q closure is added below and shows that the active quantitative task is still the row-sharp prefix atom theorem, together with the finite audit that every residual balanced-core split-pencil chart reduces to the one-parameter moving-root theorem proved here.
\end{abstract}

\maketitle




\section{The Threshold Problem}

Let \(C=\RS[\F,D,k]\) be the Reed--Solomon code obtained by evaluating polynomials of degree \(<k\) on a smooth multiplicative or circle-type domain \(D\) of size \(n\), and put \(\rho=k/n\).  In the deployed cases \(D\) is defined over a base field \(\B\subseteq\F\); write \(\beta=\log_2|\B|\).  The quantity of interest is the largest radius \(\delta\) for which the support-wise mutual correlated agreement error is at most the challenge target:
\[
        \delta_C^*(\varepsilon^*)=\sup\{\delta<1-\rho:\emca(C,\delta)\le\varepsilon^*\}.
\]
All statements below use the closed integer-ball convention: agreement at least \(a\) means radius \(1-a/n\), and an adjacent threshold certificate has an unsafe statement at \(a_0\) and a safe statement at \(a_0+1\).

The entropy--subfield candidate, used below as the certificate target, is the envelope
\[
        \delta_{\rm env}(\rho,\beta)=1-\rho-g^*(\rho,\beta),
        \qquad
        g^*(\rho,\beta)=\sup\{g:\Htwo(\rho+g)\ge \beta g\}.
\]
The point of the formula is that the base field \(\B\), not a larger challenge field, controls the frontier unless a separate transfer theorem is proved.

\begin{definition}[support-wise CA and MCA at agreement \(a\)]\label{def:ca-mca}
Let \(C\subseteq\F^D\) be linear, \(|D|=n\), and \(1\le a\le n\).  A word \(h\in\F^D\) is \(C\)-explained on \(S\subseteq D\) if some \(c\in C\) agrees with \(h\) on \(S\).  A pair \((f_1,f_2)\) is \(C^{\equiv2}\)-explained on \(S\) if some \(c_1,c_2\in C\) agree with \(f_1,f_2\) on \(S\).

For a pair \((f_1,f_2)\), a finite slope \(\gamma\in\F\) is MCA-bad at agreement \(a\) if there is a set \(S\subseteq D\), \(|S|\ge a\), such that \(f_1+\gamma f_2\) is \(C\)-explained on \(S\), but \((f_1,f_2)\) is not \(C^{\equiv2}\)-explained on that same \(S\).  It is CA-bad at agreement \(a\) if \(f_1+\gamma f_2\) is \(C\)-explained on some \(S\) with \(|S|\ge a\), while \((f_1,f_2)\) is not \(C^{\equiv2}\)-explained on any set \(T\subseteq D\) with \(|T|\ge a\).

With the uniform slope sampler \(\gamma\leftarrow\F\), define
\[
        \emca(C,1-a/n)
        =
        \frac{1}{|\F|}
        \max_{f_1,f_2\in\F^D}
        |\{\gamma:\gamma\text{ is MCA-bad at agreement }a\}|,
\]
and define \(\eca(C,1-a/n)\) analogously with CA-bad slopes.
\end{definition}

\begin{lemma}[comparison and monotonicity]\label{lem:basic-staircase}
For every linear code \(C\subseteq\F^D\) and every agreement threshold \(a\),
\[
        \eca(C,1-a/n)\le \emca(C,1-a/n).
\]
Moreover the bad-slope sets are monotone in the agreement threshold: if \(a'\le a\), then every MCA-bad slope at agreement \(a\) is MCA-bad at agreement \(a'\).  Consequently the radius function \(\delta\mapsto\emca(C,\delta)\) is nondecreasing under the closed integer-ball convention.
\end{lemma}

\begin{proof}
If a slope is CA-bad for \((f_1,f_2)\), then some support \(S\) of size at least \(a\) explains \(f_1+\gamma f_2\), and no support of size at least \(a\) explains the pair.  In particular, the same \(S\) does not explain the pair, so the slope is MCA-bad.  Taking maxima gives the first inequality.

For monotonicity, the same witness support \(S\) used at agreement \(a\) also satisfies \(|S|\ge a'\) when \(a'\le a\).  Thus it witnesses badness at the lower agreement threshold.  Since increasing the closed radius lowers the required integer agreement, \(\emca\) is nondecreasing in \(\delta\).
\end{proof}

\begin{definition}[staircase numerator]\label{def:staircase}
For an MCA row, let \(B_C^{\rm MCA}(a)\) be the maximum number of MCA-bad line parameters with a witness support of size at least \(a\).  For a list row, let \(B_C^{\rm list}(a)\) be the maximum number of codewords agreeing with a received word on at least \(a\) positions.  When the row track is fixed, write \(B(a)\) for the relevant numerator.  For target \(\varepsilon^*\) and challenge denominator \(Q\), write \(B^*=\lfloor \varepsilon^* Q\rfloor\).  An adjacent certificate at \(a_0\) is the pair of inequalities \(B(a_0)>B^*\) and \(B(a_0+1)\le B^*\).
\end{definition}

\begin{lemma}[integer budget convention]\label{lem:integer-budget}
Let a finite sampler have denominator \(Q\), and let \(B\) be the number of bad samples.  For \(B^*=\lfloor\varepsilon^*Q\rfloor\),
\[
        \frac{B}{Q}\le\varepsilon^*
        \quad\Longleftrightarrow\quad
        B\le B^*,
        \qquad
        \frac{B}{Q}>\varepsilon^*
        \quad\Longleftrightarrow\quad
        B>B^* .
\]
Thus every finite adjacent certificate can be stated purely as an integer comparison with \(B^*\).
\end{lemma}

\begin{proof}
Since \(B\) is an integer, \(B\le\varepsilon^*Q\) is equivalent to \(B\le\lfloor\varepsilon^*Q\rfloor\).  The strict inequality is the negation of the first condition.
\end{proof}

The unsafe half of the present program supplies the first inequality in \cref{def:staircase} by identity-prefix floors.  The safe half must prove the second inequality after all tangent, quotient, extension, and aperiodic residual cells have been paid without double-counting.

\begin{definition}[first-match upper ledger]\label{def:first-match-ledger}
Fix a row, an agreement \(a\), and a received line \((f_1,f_2)\).  A first-match upper ledger for this line is a finite ordered list of cells \(\mathcal C_1,\ldots,\mathcal C_s\), together with supersets \(E_i\subseteq\F\) and certified bounds \(U_i\), such that every MCA-bad slope at agreement \(a\) lies in \(\bigcup_i E_i\) and \(|E_i|\le U_i\).  The associated first-match assignment sends a bad slope \(\gamma\) to the least \(i\) with \(\gamma\in E_i\).
\end{definition}

\begin{lemma}[first-match upper ledger]\label{lem:first-match-ledger}
If every received line admits a first-match upper ledger at agreement \(a\) with \(\sum_iU_i\le U\), then \(B_C^{\rm MCA}(a)\le U\).  In particular, \(U\le B^*\) gives the MCA safe half of the adjacent certificate at agreement \(a\).
\end{lemma}

\begin{proof}
For a fixed line, the first-match sets
\[
        F_i=\{\gamma:\gamma\text{ is bad and }i\text{ is the least index with }\gamma\in E_i\}
\]
are disjoint and cover all bad slopes.  Since \(F_i\subseteq E_i\), the number of bad slopes is at most \(\sum_i|E_i|\le\sum_iU_i\le U\).  Taking the maximum over received lines gives \(B_C^{\rm MCA}(a)\le U\).
\end{proof}

\begin{remark}[MCA versus list rows]\label{rem:mca-list-route}
The four deployed rows use the same staircase language but not the same lower-side route.  The list rows compare the prefix list floor directly with \(B^*\).  The MCA rows first form a list floor for \(K=k+1\), then use the deep-point conversion to obtain a bad-slope lower bound.  In the moved v13 raw rows this conversion is packaged into an exact integer comparison against \(B^*\), but a certificate must still print which numerator, denominator, and route it uses.
\end{remark}

\begin{remark}[audit convention]\label{rem:audit-convention}
Every theorem, proposition, lemma, and corollary below is meant as a proved statement under its displayed hypotheses.  A statement is not a proof of the deployed adjacent safe row unless it explicitly supplies the upper inequality \(B(a_0+1)\le B^*\) for the row numerator in \cref{def:staircase}.  Lower-route stopping statements, quotient-rung vetoes, and raw support censuses are therefore kept separate from upper ledgers.  This convention is used to avoid hiding a dense-frontier assumption behind terminology.
\end{remark}

\section{What Is Already Banked}

The current CAP25 v13 raw manuscript records exact unsafe deployed edges.  We treat those as the active experimental source and do not reproduce its long proof here.  The table states the consequence needed by this note.

\begin{center}
\begin{tabular}{@{}llll@{}}
\toprule
row & target & proved unsafe agreement \(a_0\) & unsafe radius\\
\midrule
KoalaBear MCA & \(2^{-128}\) & \(1116047\) & \(981105/2097152\)\\
KoalaBear list & \(2^{-128}\) & \(1116046\) & \(490553/1048576\)\\
Mersenne-31 MCA & \(2^{-100}\) & \(1116023\) & \(981129/2097152\)\\
Mersenne-31 list & \(2^{-100}\) & \(1116022\) & \(490565/1048576\)\\
\bottomrule
\end{tabular}
\end{center}

Each entry is an exact integer comparison of the form
\[
        \binom{n}{m}>|\B|^w B^*
\]
after applying the route specified in \cref{rem:mca-list-route}, with the adjacent row failing in the same scanner.  The next-step bit margins for the safe inequality are approximately \(22.2,22.0,3.3,3.1\) bits in the displayed order, so a qualitative \(\poly(n)\) loss cannot settle the finite deployed rows.

\begin{center}\footnotesize
\begin{tabular}{@{}llll@{}}
\toprule
row & adjacent agreement & spare margin & current adjacent status\\
\midrule
KoalaBear MCA & \(1116048\) & \(22.1969\) bits & unsafe route stops; no \(U\)-certificate\\
KoalaBear list & \(1116047\) & \(22.0109\) bits & unsafe route stops; no \(U\)-certificate\\
Mersenne-31 MCA & \(1116024\) & \(3.2589\) bits & tight dyadic watch item, still below budget\\
Mersenne-31 list & \(1116023\) & \(3.0730\) bits & tight dyadic watch item, still below budget\\
\bottomrule
\end{tabular}
\end{center}

The exact challenge budgets are
\[
        B^*_{\mathrm{KB}}=\bigl\lfloor p_{\mathrm{KB}}^6/2^{128}\bigr\rfloor
        =274980728111395087,
        \qquad
        B^*_{\mathrm{M31}}=\bigl\lfloor p_{\mathrm{M31}}^4/2^{100}\bigr\rfloor
        =16777215.
\]
The Mersenne-31 rows therefore use target \(2^{-100}\), not \(2^{-128}\); at \(2^{-128}\) their line field would be below the target denominator and the row would be degenerate.  The current packet family records these budgets in
\[
\texttt{experimental/data/certificates/frontier-adjacent/*.packet.json}
\]
and the lightweight consistency verifier is
\[
\texttt{experimental/scripts/verify\_frontier\_adjacent\_v13\_rows.py}.
\]
That verifier checks the stored row constants, the unsafe-side replay, and the full dyadic rung-margin audit.  It does not prove \(U(a_0+1)\le B^*\), and the JSON packets correctly leave the safe certificate status unset.

\begin{remark}[finite adjacent artifact status]\label{rem:finite-artifacts}
For the MCA rows, the active lower certificate uses the quantitative deep-list conversion.  The packet values are
\[
\begin{array}{c|cc}
\text{row} & M(a_0) & M(a_0+1)\\
\hline
\text{KoalaBear MCA} & 138634741058327852652 & 57198030366\\
\text{Mersenne-31 MCA} & 4281388998575706 & 1752700 .
\end{array}
\]
Thus \(M(a_0)>B^*\) and the same lower route fails at \(a_0+1\).  This failure is not a safe upper bound.  The dyadic rung audit currently has no firing frontier rung at the adjacent agreements; the tightest item is the Mersenne-31 MCA \(c=2048\) quotient watch item, with exact mass \(12769758<16777215=B^*_{\mathrm{M31}}\).  The extension-cell artifact reduces the full-extension branch to zero-dimensional degree ceilings \(4807520,4226236,9,8\) for the four rows in table order.  These are useful target ceilings, not payments of the corresponding cells.
\end{remark}

\begin{remark}[packet convention]\label{rem:packet-convention}
The MCA rows in the current packet family contain historical adjacent-pair fields as well as the active moved-pair block.  For v13 raw, the active MCA edges are the moved pairs displayed in this note: KoalaBear MCA at \(a_0=1116047\) and Mersenne-31 MCA at \(a_0=1116023\).  The list rows are unchanged.  In the verifier terminology, the broad dyadic rung scan is gate G6 for the original packet-pair tables, while gate G7 recomputes the moved MCA lower route and the named Mersenne-31 \(c=2048\) watch item.  Thus the rung audit is a veto on currently printed quotient lower-floor refutations at \(a_0+1\), not a quotient-image upper bound.  The Mersenne-31 MCA watch item is close but quiet: it is short by \(16777215-12769758=4007457\) objects.
\end{remark}

\begin{theorem}[quantitative simple-pole list-to-MCA floor]\label{thm:simple-pole-list-floor}
Let \(C=\RS[\F,D,k]\), let \(C^+=\RS[\F,D,k+1]\), put \(q=|\F|\), \(n=|D|\), and assume \(q>n\).  Let \(U:D\to\F\) have \(L\ge1\) distinct codewords of \(C^+\) in the closed Hamming ball of radius \(\delta_0\), and assume \(\delta_0<1-k/n\).  Then for every \(\delta\in[\delta_0,1-k/n)\) there is a received line for \(C\) with at least
\[
        \left\lceil \frac{L(q-n)}{q-n+k(L-1)}\right\rceil
\]
CA-bad finite slopes at radius \(\delta\).  If \(\delta>0\), the same slopes are support-wise MCA-bad.  In particular,
\[
        \eca(C,\delta),\ \emca(C,\delta)
        \ge
        \frac{1}{q}
        \left\lceil \frac{L(q-n)}{q-n+k(L-1)}\right\rceil
\]
for \(\delta>0\), with the \(\emca\) conclusion omitted only at the degenerate endpoint \(\delta=0\).
\end{theorem}

\begin{proof}
Let \(P_1,\ldots,P_L\in\F[X]_{<k+1}\) be distinct polynomials from the \(\delta_0\)-ball around \(U\).  Since \(\delta<1-k/n\), every \(P_i\) agrees with \(U\) on more than \(k\) points.  For \(\alpha\in\Omega:=\F\setminus D\), define the simple-pole line
\[
        f_\alpha(x)=\frac{U(x)}{x-\alpha},
        \qquad
        g_\alpha(x)=\frac{-1}{x-\alpha}
        \qquad(x\in D).
\]
If \(P_i\) agrees with \(U\) on \(S_i\), then on \(S_i\)
\[
        f_\alpha(x)+P_i(\alpha)g_\alpha(x)
        =\frac{P_i(x)-P_i(\alpha)}{x-\alpha},
\]
which is the evaluation of a degree-\(<k\) polynomial.  Thus each distinct value of \(P_i(\alpha)\) gives a close slope for the line.

The pair \((f_\alpha,g_\alpha)\) is not column-close at radius \(\delta\): a common explanation on more than \(k\) points would in particular explain \(g_\alpha\) there, so \((X-\alpha)G(X)+1\) would be a polynomial of degree at most \(k\) with more than \(k\) roots and value \(1\) at \(X=\alpha\).  This is impossible.  Therefore all close slopes just constructed are CA-bad, and because the same explaining support cannot explain the pair, they are also MCA-bad whenever the radius is nonzero.

It remains to choose \(\alpha\) with many distinct values.  For \(i\ne j\), the nonzero polynomial \(P_i-P_j\) has degree at most \(k\), so \(P_i(\alpha)=P_j(\alpha)\) for at most \(k\) values \(\alpha\in\Omega\).  Hence some \(\alpha\) has at most \(k\binom L2/(q-n)\) colliding unordered pairs among the \(L\) values.  If the value multiplicities are \(m_1,\ldots,m_M\), then
\[
        \sum_r m_r=L,\qquad
        \sum_r m_r^2
        =L+2\sum_r\binom{m_r}{2}
        \le L+\frac{kL(L-1)}{q-n}.
\]
Cauchy--Schwarz gives
\[
        L^2\le M\sum_r m_r^2,
        \qquad
        M\ge \frac{L(q-n)}{q-n+k(L-1)}.
\]
Since \(M\) is an integer, the displayed ceiling is valid.
\end{proof}

\begin{corollary}[identity-prefix unsafe criterion]\label{cor:identity-prefix-unsafe}
Let \(D\subseteq\B\subseteq\F\), \(|D|=n\), and write \(b=|\B|\), \(q=|\F|\).  Fix an agreement \(a\) and a target numerator \(B^*\).
\begin{enumerate}[label=\textup{(\roman*)}]
\item For the list row \(C_K=\RS[\F,D,K]\), put \(w=a-K\).  If \(w\ge0\) and
\[
        \left\lceil\frac{\binom na}{b^w}\right\rceil>B^*,
\]
then there is a received word with more than \(B^*\) codewords of \(C_K\) agreeing on at least \(a\) positions.
\item For the MCA row \(C=\RS[\F,D,k]\), put \(K=k+1\), \(w=a-K\), and
\[
        L_a=\left\lceil\frac{\binom na}{b^w}\right\rceil .
\]
If \(w\ge0\), \(q>n\), and
\[
        \left\lceil \frac{L_a(q-n)}{q-n+k(L_a-1)}\right\rceil>B^*,
\]
then \(C\) has more than \(B^*\) support-wise MCA-bad finite slopes at agreement \(a\) for some received line.
\end{enumerate}
\end{corollary}

\begin{proof}
For (i), \cref{prop:prefix-witness} gives a prefix value whose fiber has at least \(\lceil\binom na/b^w\rceil\) members; those members are distinct degree-\(<K\) codewords agreeing with one received word on exactly \(a\) positions.

For (ii), apply the same list construction to \(C^+=\RS[\F,D,k+1]\) and then use \cref{thm:simple-pole-list-floor} at \(\delta=1-a/n\).  Since \(a\ge k+1\), this radius is \(<1-k/n\), and the displayed ceiling is a lower bound on the number of MCA-bad finite slopes.
\end{proof}

\begin{proposition}[finite packet consequences]\label{prop:finite-packet-consequences}
For the four deployed rows in the table above, let \(p\) denote the row's base prime.  The current packet family proves the unsafe inequality at the displayed active agreement \(a_0\).  More precisely:
\begin{enumerate}[label=\textup{(\roman*)}]
\item For the two list rows, the identity-prefix list floor
\[
        L(a)=\left\lceil\frac{\binom na}{p^{a-k}}\right\rceil
\]
has the exact values
\[
\begin{array}{c|cc}
\text{row} & L(a_0) & L(a_0+1)\\
\hline
\text{KoalaBear list} & 157702518233425975347 & 65065153468\\
\text{Mersenne-31 list} & 4870025984688527 & 1993678 .
\end{array}
\]
Thus \(L(a_0)>B^*\) and \(L(a_0+1)\le B^*\).
\item For the two MCA rows, with \(K=k+1\),
\[
        L(a)=\left\lceil\frac{\binom na}{p^{a-K}}\right\rceil,
        \qquad
        M(a)=\left\lceil \frac{L(a)(q-n)}{q-n+k(L(a)-1)}\right\rceil ,
\]
the active moved-pair packet values are
\[
\begin{array}{c|cc}
\text{row} & M(a_0) & M(a_0+1)\\
\hline
\text{KoalaBear MCA} & 138634741058327852652 & 57198030366\\
\text{Mersenne-31 MCA} & 4281388998575706 & 1752700 .
\end{array}
\]
Thus \(M(a_0)>B^*\) and \(M(a_0+1)\le B^*\).
\end{enumerate}
Consequently the listed \(a_0\)'s are genuine unsafe agreements for their rows, but the adjacent inequalities \(L(a_0+1)\le B^*\) and \(M(a_0+1)\le B^*\) are only lower-route failures, not safe upper bounds.
\end{proposition}

\begin{proof}
For list rows, the displayed integers are obtained by substituting the row values of \(n,k,p,a_0\) into the exact formula for \(L(a)\).  Comparing them with
\[
        B^*_{\mathrm{KB}}=274980728111395087,
        \qquad
        B^*_{\mathrm{M31}}=16777215
\]
gives \(L(a_0)>B^*\) and \(L(a_0+1)\le B^*\).  By \cref{cor:identity-prefix-unsafe}(i), the first inequality is exactly the unsafe list certificate; the second says only that this particular prefix-floor construction stops one agreement later.

For MCA rows, \cref{cor:identity-prefix-unsafe}(ii) says that the resulting bad-slope numerator is at least \(M(a)\).  Hence \(M(a_0)>B^*\) proves \(\emca(C,1-a_0/n)>\varepsilon^*\).  The displayed values are the exact integer values recorded in the active moved-pair packet fields and recomputed by the replay script.  Since the same formula gives \(M(a_0+1)\le B^*\), the same lower construction no longer proves failure at the adjacent agreement.  This is not an upper bound on all bad slopes, because other tangent, quotient, extension, Q, BC, SP, list, or sparse cells may still contribute.
\end{proof}

\begin{proposition}[dyadic rung-veto consequence]\label{prop:rung-veto}
Assume the frontier-adjacent packet verifier described above is replayed successfully.  Then no currently printed dyadic quotient lower-floor rung among the scanned scales \(c=2^j\), \(0\le j\le20\), and the printed graded-prefix, quotient-remainder, and planted-core profiles refutes any of the four active adjacent candidates \(a_0+1\).  The tightest active frontier watch item is the Mersenne-31 MCA \(c=2048\) rung with exact mass
\[
        12769758<16777215=B^*_{\mathrm{M31}}.
\]
Therefore a new adjacent unsafe certificate must either sharpen a quiet watch item, use an unscanned quotient mechanism, or come from a genuinely aperiodic, extension, sparse/CA, or list/interleaving residual branch.
\end{proposition}

\begin{proof}
This is a scope statement for the exact rung audit, not a new mathematical upper bound.  The verifier recomputes, by exact integer arithmetic, every packet field in the dyadic rung tables for the scanned scales and profiles, including whether the corresponding lower floor fires, is within the printed tightness window, and covers the frontier agreement.  The stored and recomputed active moved-pair audit has no firing frontier rung at the adjacent agreements; the only active frontier-covering sub-one-bit watch item is the displayed Mersenne-31 MCA \(c=2048\) mass.  Since that mass is strictly below \(B^*_{\mathrm{M31}}\), the scanned printed quotient lower floors do not themselves refute the adjacent safe candidates.  The conclusion is only a veto on this lower-floor family.  It does not certify a quotient-image upper ledger \(B_Q^{\rm all}(a_0+1)\).
\end{proof}

\begin{proposition}[unconditional safe-anchor gap at the active adjacent rows]\label{prop:finite-safe-anchor-gap}
Let \(a_+\) be the active adjacent agreement \(a_0+1\) for the four deployed rows, using the moved-pair convention for the two MCA rows.  The theorem-level broad safe anchors used in this note begin at agreements
\[
        A_{\rm deep}=1747627,\qquad
        A_{\rm half}=1572864,\qquad
        A_{\rm HJ}=1790031.
\]
Here \(A_{\rm deep}\) is the self-contained deep-MCA edge, \(A_{\rm half}\) is the half-distance edge using the BCIKS20 import \cite{BCIKS20}, and \(A_{\rm HJ}\) is the half-Johnson CA edge.  At the active adjacent rows the shortfalls \(A-a_+\) are:
\[
\begin{array}{c|rrr}
\text{row} & A_{\rm deep} & A_{\rm half} & A_{\rm HJ}\\ \hline
\text{KoalaBear MCA} & 631579 & 456816 & 673983\\
\text{KoalaBear list} & 631580 & 456817 & 673984\\
\text{Mersenne-31 MCA} & 631603 & 456840 & 674007\\
\text{Mersenne-31 list} & 631604 & 456841 & 674008 .
\end{array}
\]
Consequently none of these unconditional broad safe anchors supplies the adjacent upper ledger \(U(a_+)\le B^*\).  A finite adjacent certificate must use a row-specific first-match upper ledger, or a new theorem that reaches these agreements.
\end{proposition}

\begin{proof}
The three displayed safe-anchor agreements are the theorem-level thresholds printed in the frontier-adjacent packet family and recomputed by the verifier's safe-theorem applicability audit.  The active adjacent agreements are \(1116048,1116047,1116024,1116023\), respectively; for MCA these are the moved v13 raw adjacent agreements, not the historical packet-pair fields.  Subtracting these from the three anchor agreements gives the displayed table.  Since each active \(a_+\) is strictly below every listed safe-anchor agreement, the hypotheses of those safe theorems do not reach the active adjacent row.  This is only an applicability audit; it gives no upper bound for the residual cells.
\end{proof}

\begin{proposition}[full-extension cell target]\label{prop:extension-cell-target}
In the finite-adjacent packet grammar, let \(p=|\B|\) and let \(a_+=a_0+1\).  For a list row set \(w=a_0-k\); for an MCA row set \(w=a_0-k-1\), because the MCA lower route first builds a list in \(\RS[\F,D,k+1]\).  Define the prefix-normalized extension headroom
\[
        H_{\rm ext}
        :=
        \left\lfloor \frac{p^{w+1}B^*}{\binom n{a_+}}\right\rfloor .
\]
For the four deployed adjacent rows, in the order used above, exact integer arithmetic gives
\[
\begin{array}{c|c|c|c}
\text{row} & p & a_+ & H_{\rm ext}\\ \hline
\text{KoalaBear MCA} & 2^{31}-2^{24}+1 & 1116048 & 4807520\\
\text{KoalaBear list} & 2^{31}-2^{24}+1 & 1116047 & 4226236\\
\text{Mersenne-31 MCA} & 2^{31}-1 & 1116024 & 9\\
\text{Mersenne-31 list} & 2^{31}-1 & 1116023 & 8 .
\end{array}
\]
Consequently any full-orbit extension chart paid in this prefix-normalized margin by a dimension-degree term \(\Delta p^{e_Y}\), with \(\Delta\ge1\), must have \(e_Y=0\), and then must satisfy \(\Delta\le H_{\rm ext}\) in the corresponding row.
\end{proposition}

\begin{proof}
The displayed \(H_{\rm ext}\) is just the exact floor of the row's remaining prefix-normalized multiplicative margin at \(a_+\).  The packet verifier computes it from \(n=2^{21}\), \(k=2^{20}\), the row prime \(p\), the target budget \(B^*\), and the active adjacent agreement \(a_+\), using integer binomial recurrences rather than floating point.  The KoalaBear values use \(B^*=\lfloor p^6/2^{128}\rfloor\), while the Mersenne-31 values use \(B^*=\lfloor p^4/2^{100}\rfloor\).

If an extension chart contributes \(\Delta p^{e_Y}\) against this margin and \(e_Y\ge1\), then \(\Delta p^{e_Y}\ge p\).  In all four rows \(p>H_{\rm ext}\): for KoalaBear \(p>2^{30}\) while \(H_{\rm ext}<2^{23}\), and for Mersenne-31 \(p>2^{30}\) while \(H_{\rm ext}\le9\).  Thus no positive-dimensional chart can fit the printed adjacent headroom.  The zero-dimensional case contributes \(\Delta\), so fitting the same margin is exactly the inequality \(\Delta\le H_{\rm ext}\).
\end{proof}

\begin{remark}[what the extension target does not prove]\label{rem:extension-target-nonclaim}
The proposition is a target for the unpaid \(K=\F\) full-orbit branch, not a payment of that branch.  The lower minimal-field branches \(K=\B\) and \(\B<K<\F\) are routed by confinement and descent to base or lower-arity cells in the v13 raw ledger, while a genuinely full-orbit branch still needs a zero-dimensional eliminant or descended-cycle classification.  In particular, the binding Mersenne-31 list row demands total degree at most \(8\) for the full-orbit residual; a positive-dimensional family would be a new obstruction floor rather than a harmless error term.
\end{remark}

The safe side already has two broad anchors.  First, a self-contained deep-regime theorem gives MCA error at most \(n/|\F|\) below one third of the minimum distance.  Second, below one half of the minimum distance, MCA reduces to CA plus a tangent term; with the classical BCIKS half-distance CA input \cite{BCIKS20} this gives the rate-\(1/2\) safe edge \(\delta\le 1/4\).  These are not close enough to meet the unsafe edge, but they fix the outside of the remaining band.

\begin{theorem}[deep-regime MCA bound]\label{thm:gf-deep-mca}
Let \(C\subseteq\F^n\) be an \(\F\)-linear code with minimum nonzero Hamming weight \(d\), let \(q=|\F|\), and put \(r=\lfloor\delta n\rfloor\).  If
\[
        3r\le d-1,
\]
then
\[
        \eca(C,\delta)\le \frac{r+1}{q},
        \qquad
        \emca(C,\delta)\le \frac{r+1}{q}.
\]
For \(C=\RS[\F,D,k]\), this condition is \(3r\le n-k\).
\end{theorem}

\begin{proof}
Fix a pair \((f_1,f_2)\), and let
\[
        S_{\rm cl}=\{\gamma\in\F:\dist(f_1+\gamma f_2,C)\le r\}.
\]
If \(|S_{\rm cl}|\le r+1\), then this pair contributes at most \(r+1\) close slopes, hence at most \(r+1\) CA- or MCA-bad slopes.

Assume \(|S_{\rm cl}|\ge2\).  Choose distinct \(\gamma_1,\gamma_2\in S_{\rm cl}\), codewords \(p_1,p_2\), and error sets \(E_1,E_2\), \(|E_i|\le r\), such that \(f_1+\gamma_i f_2=p_i\) off \(E_i\).  Define
\[
        c_2=\frac{p_1-p_2}{\gamma_1-\gamma_2},
        \qquad
        c_1=p_1-\gamma_1c_2,
        \qquad
        e_i=f_i-c_i .
\]
Off \(T=E_1\cup E_2\), \(|T|\le2r\), the two linear equations imply \(e_1=e_2=0\).  Let \(T'\subseteq T\) be the support of the error pair \((e_1,e_2)\).

For any \(\gamma\in S_{\rm cl}\), with explaining codeword \(p\) and error set \(E\), the codeword \(p-(c_1+\gamma c_2)\) vanishes outside \(T\cup E\), hence has weight at most \(3r\le d-1\).  It is therefore zero, so
\[
        \wt(e_1+\gamma e_2)\le r
        \qquad(\gamma\in S_{\rm cl}).
\]
If \(|T'|\ge r+1\), each \(\gamma\in S_{\rm cl}\) must make at least \(|T'|-r\) of the nonzero affine coordinate functions \(e_{1,j}+\gamma e_{2,j}\) vanish.  Each coordinate vanishes for at most one \(\gamma\).  Double counting gives
\[
        |S_{\rm cl}|(|T'|-r)\le |T'|,
        \qquad\text{hence}\qquad
        |S_{\rm cl}|\le \frac{|T'|}{|T'|-r}\le r+1.
\]
Thus this pair again contributes at most \(r+1\) close slopes.

It remains to consider \(|T'|\le r\).  Then \(f_i=c_i\) outside \(T'\), so the pair is column-close to \(C^{\equiv2}\).  It has no CA-bad slope.  For MCA, if \(\gamma\) has a witness support \(S\) of size at least \(n-r\), the explaining codeword for \(f_1+\gamma f_2\) agrees with \(c_1+\gamma c_2\) on at least \(n-2r\) positions; since \(2r\le d-1\), the explaining codeword is \(c_1+\gamma c_2\).  A simultaneous explanation of the pair on \(S\) would also be forced to be \((c_1,c_2)\).  Hence badness can occur only if \(e_1+\gamma e_2\) vanishes at some coordinate of \(T'\), and each coordinate contributes at most one slope.  There are at most \(|T'|\le r\) such slopes.  Taking the maximum over pairs proves both bounds.  For Reed--Solomon, \(d=n-k+1\).
\end{proof}

\begin{proposition}[exact high-agreement tangent cell]\label{prop:exact-tangent-cell}
Let \(C=\RS[\F,D,k]\), \(|D|=n\), assume \(0<k<n\), and let \(A=n-r\).  If \(3r\le n-k\), then the maximum, over received lines, of the number of finite MCA-bad slopes at agreement \(A\) is exactly \(r+1=n-A+1\).  The same exact value holds for CA-bad slopes.  Thus the tangent/common-line paid cell is pointwise exact throughout this range.
\end{proposition}

\begin{proof}
The upper bound is \cref{thm:gf-deep-mca} multiplied by \(|\F|\).

For the lower bound, choose \(T=\{t_0,\ldots,t_r\}\subseteq D\) and distinct \(\gamma_0,\ldots,\gamma_r\in\F\).  This is possible because \(|\F|\ge n\ge r+1\).  Define
\[
        f_2(t_i)=1,\qquad f_1(t_i)=-\gamma_i,
        \qquad
        f_1(x)=f_2(x)=0\quad(x\in D\setminus T).
\]
For slope \(\gamma_i\), the word \(f_1+\gamma_i f_2\) vanishes on
\[
        S_i=D\setminus(T\setminus\{t_i\}),
        \qquad |S_i|=n-r=A,
\]
so it is explained there by the zero codeword.

We need two elementary inequalities.  If \(r=0\), then \(n-r-1\ge k\) and \(n-2r-1\ge k\).  If \(r\ge1\), then \(r+1\le3r\le n-k\), so \(n-r-1\ge k\); also \(3r\le n-k\) gives \(2r\le n-k-1\), hence \(n-2r-1\ge k\).

First fix \(i\).  If two degree-\(<k\) polynomials explained \(f_1\) and \(f_2\) simultaneously on \(S_i\), then both would vanish on \(D\setminus T\), a set of size \(n-r-1\ge k\), hence both would be zero.  This contradicts \(f_2(t_i)=1\), since \(t_i\in S_i\).  Therefore \(\gamma_i\) is MCA-bad at agreement \(A\), and the \(r+1\) slopes are distinct.

For CA-badness, let \(S\subseteq D\) have \(|S|\ge A\).  Since \(|T|=r+1\), the set \(S\cap T\) is nonempty.  Also
\[
        |S\cap(D\setminus T)|\ge |S|-|T|\ge n-2r-1\ge k .
\]
Any simultaneous degree-\(<k\) explanation of \(f_1,f_2\) on \(S\) would therefore vanish on at least \(k\) points outside \(T\), hence be the zero pair, contradicting \(f_2(t)=1\) for \(t\in S\cap T\).  Thus the pair is CA-far at agreement \(A\), while each \(f_1+\gamma_i f_2\) is explained on \(S_i\).  The same \(r+1\) slopes are CA-bad.
\end{proof}

\begin{remark}[endpoint audit for the exact tangent cell]\label{rem:tangent-endpoint-audit}
The extra hypothesis \(k<n\) is only an endpoint issue for the lower construction when \(r=0\).  If \(k=n\), then \(C=\F^D\) and there are no bad slopes, so the exact lower value \(r+1=1\) would be false.  The upper bound in \cref{thm:gf-deep-mca} remains valid for the full-space code.
\end{remark}

\begin{theorem}[MCA from CA below half distance]\label{thm:gf-mca-from-ca}
Let \(C\subseteq\F^n\) be linear with minimum nonzero Hamming weight \(d\), \(q=|\F|\), and \(r=\lfloor\delta n\rfloor\).  If \(2r\le d-1\), then
\[
        \emca(C,\delta)\le \max\!\left(\eca(C,\delta),\frac{r}{q}\right).
\]
For \(C=\RS[\F,D,k]\), the hypothesis is \(2r\le n-k\).
\end{theorem}

\begin{proof}
Fix \((f_1,f_2)\).  If the pair is column-far from \(C^{\equiv2}\) at radius \(\delta\), then every MCA-bad slope is already CA-bad: the point \(f_1+\gamma f_2\) is close to \(C\) on its witness support, while the pair is far on every support of the required size.

If the pair is column-close, choose codewords \(p_1,p_2\in C\) agreeing with \(f_1,f_2\) outside an error set \(E\), \(|E|\le r\), and write \(e_i=f_i-p_i\).  Let \(\gamma\) be MCA-bad with witness support \(S\), \(|S|\ge n-r\), and explaining codeword \(c\).  On \(S\setminus E\), the codeword \(c-(p_1+\gamma p_2)\) vanishes; this set has size at least \(n-2r\), so the codeword has weight at most \(2r\le d-1\) and is zero.  Thus on \(S\), the error word \(e_1+\gamma e_2\) vanishes.  The support \(S\) must meet \(E\), otherwise \((p_1,p_2)\) would explain the pair on \(S\).  At some \(j\in S\cap E\), we have \(e_{1,j}+\gamma e_{2,j}=0\) with \((e_{1,j},e_{2,j})\ne(0,0)\), forcing \(e_{2,j}\ne0\) and \(\gamma=-e_{1,j}/e_{2,j}\).  Hence a column-close pair contributes at most \(|E|\le r\) MCA-bad slopes.  Taking the maximum over pairs gives the displayed inequality.
\end{proof}

\begin{theorem}[subfield confinement of bad slopes]\label{thm:subfield-confinement}
Let \(\B\subseteq\F\) be finite fields, let \(D\subseteq\B\), and let \(C=\RS[\F,D,k]\).  If \(f,g\in\B^D\), then every support-wise CA- or MCA-bad slope for the line \(f+\gamma g\) lies in \(\B\).  More precisely, if \(\gamma\in\F\setminus\B\) and \(f+\gamma g\) is explained by a codeword of \(C\) on a support \(S\subseteq D\), then the pair \((f,g)\) is simultaneously explained by two codewords of \(C\) on the same support \(S\).  Consequently every \(\B\)-valued pair has bad-slope density at most \(|\B|/|\F|\) under the uniform \(\gamma\leftarrow\F\) sampler.  The same conclusion holds after multiplying both words by a common scalar in \(\F^\times\).
\end{theorem}

\begin{proof}
Choose a \(\B\)-basis \(1=\beta_0,\beta_1,\ldots,\beta_{s-1}\) of \(\F\).  Every word \(h\in\F^D\) decomposes uniquely as \(h=\sum_i\beta_i h_i\) with \(h_i\in\B^D\).  Since \(D\subseteq\B\), every codeword of \(C=\RS[\F,D,k]\) decomposes componentwise into codewords of \(\RS[\B,D,k]\): decompose the coefficients of its defining polynomial in the same basis and evaluate on \(D\).

Let \(\gamma=\sum_i\gamma_i\beta_i\notin\B\), and choose \(i_*>0\) with \(\gamma_{i_*}\ne0\).  Suppose a codeword \(c=\sum_i\beta_i c_i\in C\) agrees with \(f+\gamma g\) on \(S\).  Comparing basis components on \(S\) gives
\[
        c_0=f+\gamma_0 g,\qquad
        c_{i_*}=\gamma_{i_*}g .
\]
Thus \(g\) agrees on \(S\) with \(\gamma_{i_*}^{-1}c_{i_*}\in C\), and \(f\) agrees on \(S\) with \(c_0-\gamma_0\gamma_{i_*}^{-1}c_{i_*}\in C\).  Hence every support explaining the line point also explains the pair on the same support.

For MCA this directly rules out badness of every \(\gamma\notin\B\).  For CA, if \(f+\gamma g\) is close on some support of size required by the field radius, the same support gives a column explanation of \((f,g)\); hence \(\gamma\notin\B\) cannot be CA-bad whenever the internal radius is at least the field radius.  Therefore all bad slopes lie in \(\B\), giving density at most \(|\B|/|\F|\).  Multiplication by a common scalar preserves all support explanations by linearity of \(C\), so the same statement holds for scalar multiples of \(\B\)-valued pairs.
\end{proof}

\section{Proved Prefix and Split-Pencil Reductions}

We record the part of the Q, BC, and SP program that is already theorem-level.  For a finite set \(D\subseteq\F\) and \(S\subseteq D\), put
\[
        \ell_S(X)=\prod_{x\in S}(X-x).
\]
For \(1\le w<m\le n\), define \(\Phi_w(M)\) to be the vector of the first \(w\) coefficients immediately below the leading term of the monic degree-\(m\) polynomial \(\ell_M\).  Thus two \(m\)-subsets have the same \(\Phi_w\)-value exactly when their locator polynomials have the same monic term and same next \(w\) high-degree coefficients.  Write \(\Fib_w(z)=\{M\subseteq D: |M|=m,\ \Phi_w(M)=z\}\).

\begin{proposition}[Newton dictionary]\label{prop:newton}
If \(\operatorname{char}\B>w\), then \(\Phi_w\) is triangularly equivalent to the power-sum map
\[
        M\longmapsto \left(\sum_{x\in M}x,\sum_{x\in M}x^2,\ldots,\sum_{x\in M}x^w\right).
\]
In particular the two maps have exactly the same fiber sizes.
\end{proposition}

\begin{proof}
Newton's identities express the elementary symmetric functions \(e_1,\ldots,e_w\) and the power sums \(p_1,\ldots,p_w\) by mutually inverse triangular polynomial changes of coordinates, because every integer \(1,\ldots,w\) is invertible in \(\B\).  The sign convention in \(\Phi_w\) is coordinatewise invertible, so the fibers are unchanged.
\end{proof}

\begin{proposition}[identity-prefix witness lists]\label{prop:prefix-witness}
Let \(K=m-w\).  For every prefix value \(z=(z_1,\ldots,z_w)\), define the received word
\[
        U_z(X)=X^m+\sum_{j=1}^{w}z_jX^{m-j}
\]
by evaluation on \(D\).  The degree-\(<K\) codewords agreeing with \(U_z\) on at least \(m\) points are exactly
\[
        U_z-\ell_M,\qquad M\in\Fib_w(z).
\]
No such codeword agrees with \(U_z\) on more than \(m\) points.  Consequently some \(z\) gives a list of size at least \(\lceil \binom{n}{m}|\B|^{-w}\rceil\).
\end{proposition}

\begin{proof}
If \(M\in\Fib_w(z)\), then the leading \(w+1\) coefficients of \(U_z\) and \(\ell_M\) agree, so \(\deg(U_z-\ell_M)<m-w=K\), and this codeword agrees with \(U_z\) on \(M\).  Conversely, if a degree-\(<K\) polynomial \(c\) agrees with \(U_z\) on at least \(m\) points, then \(U_z-c\) is monic of degree \(m\) and has those points as roots.  It must equal \(\ell_M\) for the corresponding \(m\)-set \(M\), and comparison of the first \(w\) coefficients puts \(M\) in \(\Fib_w(z)\).  A nonzero degree-\(m\) polynomial cannot vanish on more than \(m\) points.  The average fiber size over the \(|\B|^w\) possible prefixes gives the last sentence.
\end{proof}

\begin{proposition}[pole-line transport to MCA]\label{prop:pole-line}
Let \(C=\RS[\F,D,k]\), fix \(k+1\le m\le n\), and put \(K=k+1\), \(w=m-K\).  Assume \(z\in\B^w\), \(\alpha\in\F\setminus D\), and define
\[
        f_\alpha(x)=\frac{U_z(x)}{x-\alpha},
        \qquad
        g_\alpha(x)=\frac{-1}{x-\alpha}\qquad(x\in D).
\]
For a slope \(\zeta\in\F\) and a set \(S\subseteq D\) with \(|S|\ge m\), the word \(f_\alpha+\zeta g_\alpha\) is explained by \(C\) on \(S\) if and only if
\[
        |S|=m,\qquad S\in\Fib_w(z),\qquad
        \ell_S(\alpha)=U_z(\alpha)-\zeta .
\]
Moreover the pair \((f_\alpha,g_\alpha)\) is never simultaneously explained by two codewords of \(C\) on any set of size at least \(k+1\).  Hence the finite MCA-bad slopes at agreement threshold \(m\) are exactly
\[
        \bigl\{\,U_z(\alpha)-\ell_S(\alpha):S\in\Fib_w(z)\,\bigr\},
\]
up to collisions in the displayed evaluation map.
\end{proposition}

\begin{proof}
Suppose \(c=P|_D\), \(\deg P<k\), explains \(f_\alpha+\zeta g_\alpha\) on \(S\).  Then \(U_z(x)-\zeta=(x-\alpha)P(x)\) on \(S\).  The polynomial
\[
        \Psi(X)=U_z(X)-\zeta-(X-\alpha)P(X)
\]
is monic of degree \(m\), because \(\deg((X-\alpha)P)\le k\le m-1\), and it vanishes on \(S\).  Thus \(|S|=m\), \(\Psi=\ell_S\), the first \(w\) coefficients of \(\ell_S\) are those of \(U_z\), and \(\ell_S(\alpha)=U_z(\alpha)-\zeta\).  Conversely, if these three conditions hold, then \(U_z-\ell_S\) has degree at most \(k\) and takes the value \(\zeta\) at \(\alpha\), so
\[
        P(X)=\frac{U_z(X)-\ell_S(X)-\zeta}{X-\alpha}
\]
has degree \(<k\) and explains \(f_\alpha+\zeta g_\alpha\) on \(S\).

For the far condition, if \(g_\alpha\) agreed on some \(S\), \(|S|\ge k+1\), with a polynomial \(G\) of degree \(<k\), then \((X-\alpha)G(X)+1\) would have degree at most \(k\) and at least \(k+1\) roots, hence would be zero.  Its value at \(X=\alpha\) is \(1\), a contradiction.  A simultaneous explanation of the pair would in particular explain \(g_\alpha\).  The final claim follows from the support-wise MCA definition and the first part of the proof.
\end{proof}

\begin{theorem}[fiber-to-slope conversion]\label{thm:fiber-to-slope}
Let \(D\) be a multiplicative coset of order \(n\), let \(C=\RS[\F,D,k]\), and let \(k+1\le m\le n\), \(w=m-k-1\).  Fix a prefix value \(z\) and a subfamily \(\mathcal G\subseteq\Fib_w(z)\) of size \(N_0\).  Then there is a pole \(\alpha\in\F\setminus D\) for which the line \((f_\alpha,g_\alpha)\) of \cref{prop:pole-line} has at least
\[
        N_0-\frac{\binom{N_0}{2}k}{|\F|-n}
\]
distinct finite MCA-bad slopes at agreement threshold \(m\).  If \(N_0\le (|\F|-n)/k+1\), this is at least
\[
        \frac{N_0}{2}\left(1-\frac{k}{|\F|-n}\right).
\]
If, in addition, \(\gcd(m,n)=1\), every threshold-\(m\) witness support of every slope produced by this construction is aperiodic.
\end{theorem}

\begin{proof}
For distinct \(S,S'\in\mathcal G\), the polynomial \(\ell_S-\ell_{S'}\) has degree at most \(m-w-1=k\), because the leading term and the first \(w\) prefix coefficients cancel.  Hence the equality \(\ell_S(\alpha)=\ell_{S'}(\alpha)\) holds for at most \(k\) values of \(\alpha\).  Averaging over \(\Omega=\F\setminus D\), some \(\alpha\in\Omega\) has at most \(\binom{N_0}{2}k/(|\F|-n)\) colliding pairs among the values \(\ell_S(\alpha)\), \(S\in\mathcal G\).  The image size is at least the number of points minus the number of colliding pairs, giving the first bound.  The second bound is the displayed inequality with \(N_0-1\le (|\F|-n)/k\).

By \cref{prop:pole-line}, each distinct value \(U_z(\alpha)-\ell_S(\alpha)\) is an MCA-bad slope, and every threshold-\(m\) witness lies in \(\Fib_w(z)\).  Finally, if a support \(S\subseteq D\) is periodic under a nontrivial subgroup of the order-\(n\) stabilizer, then it is a union of orbits of some size \(c>1\) with \(c\mid n\), so \(c\mid |S|=m\).  When \(\gcd(m,n)=1\), this is impossible; all such supports are aperiodic.
\end{proof}

\begin{theorem}[no uniform polynomial aperiodic bound below the envelope]\label{thm:no-old-band}
There is an infinite family of smooth multiplicative rate-\(1/2\) rows for which, at agreements below the entropy--subfield envelope, a single received line has \(2^{\Omega(n)}\) MCA-bad finite slopes, all witnessed aperiodically.  Consequently no \(q\)-independent polynomial bound can hold for the unnormalized aperiodic band below the envelope.
\end{theorem}

\begin{proof}
Let \(n=2^\nu\), \(k=n/2\), and choose a prime \(p_\nu\equiv1\pmod n\) with \(p_\nu=n^{O(1)}\), available by Linnik's theorem \cite{Linnik44}.  Put \(\B_\nu=\F_{p_\nu}\), let \(D_\nu=\mu_n\subseteq\B_\nu^\times\), and write \(\beta_\nu=\log_2p_\nu=O(\log n)\).  Choose an odd \(m_\nu\) with
\[
        \frac{n}{8\beta_\nu}\le m_\nu-k\le \frac{n}{4\beta_\nu}
\]
and set \(w_\nu=m_\nu-k-1\).  Let \(\F_\nu/\B_\nu\) be an extension with \(q_\nu=|\F_\nu|\ge k\,2^{n/8}\), chosen minimally up to the next extension degree.

By \cref{prop:prefix-witness}, some prefix fiber has size at least \(\binom n{m_\nu}p_\nu^{-w_\nu}\).  Since \(g=(m_\nu-k)/n\le1/(4\beta_\nu)\), the elementary estimate \(H_2(1/2+g)\ge1-4g^2\) gives
\[
        \log_2\binom n{m_\nu}-w_\nu\beta_\nu
        \ge n\bigl(1-4g^2\bigr)-ng\beta_\nu-O(\log n)
        \ge \frac{n}{2}
\]
for all large \(\nu\).  Thus the heaviest fiber has at least \(2^{n/2}\) members.  By the minimal choice of the extension degree, \(\lceil(q_\nu-n)/k\rceil\le p_\nu 2^{n/8}=2^{n/8+O(\log n)}\), so we may choose a subfamily \(\mathcal G\) of size \(N_0=\lceil(q_\nu-n)/k\rceil\).

\Cref{thm:fiber-to-slope} gives a pole line with at least
\[
        \frac{N_0}{2}\left(1-\frac{k}{q_\nu-n}\right)\ge 2^{n/8-2}
\]
bad slopes for all large \(\nu\).  Because \(m_\nu\) is odd and \(n\) is a power of two, \(\gcd(m_\nu,n)=1\), so every witness support produced by the theorem is aperiodic.  This is exponential in \(n\), contradicting any uniform polynomial aperiodic bound in the sub-envelope band.
\end{proof}

\begin{remark}[what the previous theorem changes]\label{rem:normalized-band}
\Cref{thm:no-old-band} does not weaken the resolution program; it localizes it.  The identity-prefix floor is the expected unsafe side below the entropy--subfield envelope, so the positive aperiodic input must be normalized with its left edge at that envelope.  Q, BC, and SP are therefore used below as names for safe-side certificate cells after the identity-prefix floor, quotient pullbacks, and planted strata have already been paid.
\end{remark}

\begin{proposition}[prefix-collision rigidity]\label{prop:prefix-rigidity}
If \(M\ne M'\) are \(m\)-subsets of \(D\) with \(\Phi_w(M)=\Phi_w(M')\), then
\[
        |M\setminus M'|=|M'\setminus M|\ge w+1.
\]
If equality holds, then \(\ell_{M\setminus M'}-\ell_{M'\setminus M}\) is a nonzero constant.
\end{proposition}

\begin{proof}
Let \(R=M\cap M'\), put \(A=\ell_{M\setminus M'}\), \(B=\ell_{M'\setminus M}\), and \(G=\ell_R\).  If \(e=|M\setminus M'|\), then
\[
        \ell_M-\ell_{M'}=G(A-B).
\]
The equality of the first \(w\) high-degree coefficients gives \(\deg(\ell_M-\ell_{M'})\le m-w-1\), while \(\deg G=m-e\).  Hence \(\deg(A-B)\le e-w-1\).  If \(e\le w\), this forces \(A=B\), and then the two disjoint root sets \(M\setminus M'\) and \(M'\setminus M\) coincide; so \(e=0\), contrary to \(M\ne M'\).  Thus \(e\ge w+1\).  When \(e=w+1\), the same inequality gives \(\deg(A-B)\le0\); since the root sets are disjoint and nonempty, \(A\ne B\), so the constant is nonzero.
\end{proof}

\begin{proposition}[exact second moment]\label{prop:second-moment}
Let \(N_w(z)=|\Fib_w(z)|\).  For \(D'\subseteq D\), let \(\operatorname{sp}_w(e;D')\) be the number of ordered pairs \((A,B)\) of monic degree-\(e\) polynomials split over disjoint root sets in \(D'\) such that \(\deg(A-B)\le e-w-1\).  Then
\[
        \sum_z N_w(z)^2
        =
        \binom{n}{m}
        +
        \sum_{e=w+1}^{\min(m,n-m)}
        \sum_{\substack{R\subseteq D\\ |R|=m-e}}
        \operatorname{sp}_w(e;D\setminus R).
\]
The top nonzero stratum \(e=w+1\) consists exactly of constant-shift split pairs \(A\) and \(A-c\), \(c\in\F^\times\), both split over \(D'\).
\end{proposition}

\begin{proof}
The left side counts ordered pairs \((M,M')\) of \(m\)-subsets with the same prefix.  The diagonal contributes \(\binom{n}{m}\).  For \(M\ne M'\), set \(R=M\cap M'\), \(A=\ell_{M\setminus M'}\), \(B=\ell_{M'\setminus M}\), and \(e=|M\setminus M'|\).  This is a bijective encoding of off-diagonal pairs by \(R\) together with a pair \((A,B)\) split over \(D\setminus R\) with disjoint root sets.  The same degree computation as in \cref{prop:prefix-rigidity} shows that prefix equality is precisely \(\deg(A-B)\le e-w-1\).  The range \(e\ge w+1\) follows from \cref{prop:prefix-rigidity}, and \(e\le\min(m,n-m)\) is forced by the sizes of \(M\) and \(D\setminus R\).  If \(e=w+1\), then \(\deg(A-B)\le0\), so \(A-B=c\in\F^\times\), and the converse is immediate.
\end{proof}

\begin{corollary}[anticode fiber cap]\label{cor:anticode-cap}
For every \(z\), one has
\[
        |\Fib_w(z)|\le \frac{\binom{n}{m-w}}{\binom{m}{w}}.
\]
\end{corollary}

\begin{proof}
By \cref{prop:prefix-rigidity}, two distinct fiber members cannot contain the same \((m-w)\)-subset.  Counting pairs \((R,M)\) with \(M\in\Fib_w(z)\), \(R\subseteq M\), and \(|R|=m-w\) gives the claim.
\end{proof}

\begin{proposition}[moment sandwich]\label{prop:moment-sandwich}
Let \(\mu(z)=|\Fib_w(z)|/\binom{n}{m}\), set
\[
        \Gamma_r=|\B|^{w(r-1)}\sum_z\mu(z)^r,\qquad r\ge2,
\]
and let \(R=|\B|^w\max_z\mu(z)\).  Then
\[
        \frac{\log_2\Gamma_r}{r-1}\le \log_2 R
        \le \frac{\log_2\Gamma_r+w\log_2|\B|}{r}.
\]
\end{proposition}

\begin{proof}
Since \(\mu(z)^r\le(\max_z\mu(z))^{r-1}\mu(z)\), we have \(\Gamma_r\le R^{r-1}\).  Conversely \(\max_z\mu(z)\le(\sum_z\mu(z)^r)^{1/r}\), and multiplying by \(|\B|^w\) gives the upper bound for \(R\).
\end{proof}

\begin{theorem}[moment criterion for Q]\label{thm:moment-q}
With the notation of \cref{prop:moment-sandwich}, suppose a quotient/planted-pruned prefix problem has normalized maximum-fiber ratio
\[
        R=|\B|^w\max_z\frac{|\Fib_w(z)|}{\binom nm}.
\]
Let \(r\ge2\).  If
\[
        \Gamma_r\le \exp(A)
        \qquad\text{and}\qquad
        w\log|\B|\le B,
\]
then
\[
        R\le \exp\!\left(\frac{A+B}{r}\right).
\]
Consequently, an asymptotic Q theorem follows from any sequence of moment bounds with \(A=o(rn)\) and \(B/r=o(n)\).  For a finite adjacent row with remaining bit margin \(\Delta>0\), this moment route can certify \(R\le2^\Delta\) only if
\[
        \log_2\Gamma_r+w\log_2|\B|\le r\Delta .
\]
\end{theorem}

\begin{proof}
The first bound is the upper half of \cref{prop:moment-sandwich}, rewritten with natural logarithms.  If \(A=o(rn)\) and \(B/r=o(n)\), then \(\log R=o(n)\), which is the \(e^{o(n)}\) form of Q after the paid quotient and planted strata have been removed.

For the finite statement, take logarithms base \(2\) in the same inequality:
\[
        \log_2R\le\frac{\log_2\Gamma_r+w\log_2|\B|}{r}.
\]
To prove \(R\le2^\Delta\), the numerator must be at most \(r\Delta\).
\end{proof}

\begin{definition}[primitive logarithmic collision moment]\label{def:primitive-logmoment}
After all first-match paid cells have been removed, let \(N^{\rm prim}(s)\) be the number of remaining primitive prefix-boundary supports with prefix value \(s\).  Put
\[
        \overline N=\binom nm|\B|^{-w},
        \qquad
        R^{\rm prim}(s)=\frac{N^{\rm prim}(s)}{\overline N},
        \qquad
        M_{\rm prim}=\max_sR^{\rm prim}(s),
\]
and for an integer \(r\ge2\),
\[
        \Gamma_r^{\rm prim}
        =
        |\B|^{-w}\sum_s R^{\rm prim}(s)^r .
\]
\end{definition}

\begin{theorem}[logarithmic moments are equivalent to primitive asymptotic Q]\label{thm:logmoment-equivalence}
Let \(r=r(n)\to\infty\) and assume
\[
        \frac{w\log|\B|}{r}=o(n).
\]
Then the following two asymptotic statements are equivalent:
\[
        \Gamma_r^{\rm prim}\le\exp(o(n)r)
        \qquad\Longleftrightarrow\qquad
        M_{\rm prim}\le\exp(o(n)).
\]
Thus primitive asymptotic Q is exactly a logarithmic-order collision-moment theorem.  For fixed-bit base fields, \(r=\lceil\log n\rceil\) is enough whenever \(w=O(n)\).  If \(\log|\B|=O(\log n)\), one may take \(r=(\log n)^2\); more generally any \(r\to\infty\) with \(w\log|\B|/r=o(n)\) suffices.
\end{theorem}

\begin{proof}
First suppose \(\Gamma_r^{\rm prim}\le\exp(o(n)r)\).  Since
\[
        M_{\rm prim}^r
        \le
        \sum_sR^{\rm prim}(s)^r
        =
        |\B|^w\Gamma_r^{\rm prim},
\]
we have
\[
        \log M_{\rm prim}
        \le
        \frac{w\log|\B|+\log\Gamma_r^{\rm prim}}{r}
        =
        o(n).
\]
This proves \(M_{\rm prim}\le\exp(o(n))\).

Conversely, suppose \(M_{\rm prim}\le\exp(o(n))\).  Since the primitive family is a subfamily of all \(m\)-subsets,
\[
        \sum_sR^{\rm prim}(s)
        =
        \frac{\sum_sN^{\rm prim}(s)}{\binom nm|\B|^{-w}}
        \le
        |\B|^w .
\]
Therefore
\[
        \Gamma_r^{\rm prim}
        =
        |\B|^{-w}\sum_sR^{\rm prim}(s)^r
        \le
        M_{\rm prim}^{r-1}
        \le
        \exp(o(n)r),
\]
as required.
\end{proof}

\begin{problem}[primitive entropic inverse theorem for Vandermonde slice sums]\label{prob:entropy-inverse-q}
Prove the following standalone additive-combinatorics statement, or identify the extra obstruction cell that must be added to the first-match ledger.  Let \(\K=\K_N\) be a finite field, let \(T=T_N\subseteq\K\) have \(|T|=N\), let \(R=R_N\asymp N\), and choose nonzero weights \(\rho(t)\in\K^\times\).  Define the Vandermonde columns
\[
        v_t=\rho(t)(1,t,t^2,\ldots,t^{R-1})\in\K^R,
        \qquad t\in T .
\]
Let \(\Omega\subseteq\{0,1\}^T\), or more generally \(\Omega\subseteq\{-1,0,1\}^T\), be a fixed-density profile slice: all active coordinate counts lie between \(\alpha N\) and \((1-\alpha)N\) for some fixed \(\alpha>0\).  Let
\[
        \Phi:\Omega\to\K^R,
        \qquad
        \Phi(x)=\sum_{t\in T}x_tv_t,
\]
and let \(\Omega^\circ\subseteq\Omega\) be the primitive residual model after quotient pullbacks, Chebyshev/dihedral pullbacks, planted common blocks, tangent cells, extension cells, differential-locator low-defect cells, and saturation cells have been removed.  Assume the frontier normalization
\[
        \log|\Omega^\circ|-R\log|\K|=o(N).
\]
For \(X_0\) uniform on \(\Omega^\circ\), write \(\nu=\Phi_*X_0\), and define the normalized Renyi collision excess
\[
        \Gamma_\ell(\Omega^\circ,\Phi)
        =
        |\K|^{R(\ell-1)}\sum_{s\in\K^R}\nu(s)^\ell
        =
        |\K|^{-R}\sum_s
        \left(
        \frac{|\Omega^\circ\cap\Phi^{-1}(s)|}
             {|\Omega^\circ|\,|\K|^{-R}}
        \right)^\ell .
\]
For every fixed \(\alpha,\kappa,\eta>0\) and \(A<\infty\), with \(\kappa N\le R\le\kappa^{-1}N\), if \(\ell=\ell_N\to\infty\), \(\ell\le(\log N)^A\), and
\[
        \Gamma_\ell(\Omega^\circ,\Phi)\ge\exp(\eta N\ell),
\]
then at least one of the following alternatives holds.

\begin{enumerate}[label=\textup{(\alph*)}]
\item A positive-density restriction of the datum lies in one of the bounded-complexity algebraic cells already removed above.
\item There is a positive-density set \(U\subseteq T\) for which
\[
        \operatorname{rank}_{\K}\Span\{v_t:t\in U\}<\min(|U|,R).
\]
\end{enumerate}

The robust proof form may pass through an additional Sidon/free-energy branch: either level-typical star trades have entropy-small difference \(H(Y-Y')\le H(Y)+o(N)\), leading by entropic BSG/PFR and a slice-derivative lemma to the rank-defect alternative above, or else the free-energy decay of \(Y-Y'\) bounds that level's contribution to \(\Gamma_\ell\) by \(\exp(-\Omega(N\ell))\).  This problem is intentionally stated without RS, decoding, MCA, Q, or SP language.  It is the independent inverse theorem needed by the primitive Q route.
\end{problem}

\begin{remark}[why standard inverse theorems do not plug in directly]\label{rem:standard-inverse-gap}
The intended proof route follows the Tao--Vu inverse Littlewood--Offord and Green--Ruzsa/Freiman template \cite{TaoVu09,GreenRuzsa07}; the bounded-torsion PFR theorem of Gowers--Green--Manners--Tao \cite{GGMT24} is also relevant once an entropy-small-doubling object has been extracted.  The missing point is the scale.  Standard inverse Littlewood--Offord theorems are usually triggered by polynomial concentration, while the random prefix-fiber scale here is \(|\B|^{-w}=\exp(-\Theta(n))\).  A Q failure may be only \(\exp(\eta n)|\B|^{-w}\), still exponentially small.  The needed input is therefore a Renyi-entropy or logarithmic-moment inverse theorem for primitive moment-curve subset sums, not a black-box invocation of the polynomial-concentration inverse theorem.
\end{remark}

\begin{proposition}[Vandermonde rank kills the inverse output]\label{prop:vandermonde-kills-low-rank}
Let \(y_1,\ldots,y_t\) be distinct elements of \(T\), with \(t\le R\) and all weights \(\rho(y_i)\ne0\).  Then the columns
\[
        v_{y_i}=(\rho(y_i),\rho(y_i)y_i,\ldots,\rho(y_i)y_i^{R-1})
\]
are linearly independent.  Consequently, if \(R=\Theta(n)\), no affine subspace of rank \(o(n)\) contains a positive-density subset of the primitive moment-curve columns.
\end{proposition}

\begin{proof}
Multiply the column \(v_{y_i}\) by \(\rho(y_i)^{-1}\).  The resulting \(t\times t\) leading minor is the Vandermonde matrix \((y_i^{j})_{0\le j<t}\), whose determinant is \(\prod_{i<j}(y_j-y_i)\ne0\).  Hence any \(t\le R\) distinct columns are independent.  A set of \(cn\) columns has affine span of dimension at least \(cn-1\) as long as \(cn\le R\), which is incompatible with affine rank \(o(n)\) for fixed \(c>0\) and large \(n\).
\end{proof}

\begin{corollary}[asymptotic Q from the entropy inverse theorem]\label{cor:asymp-q-from-entropy-inverse}
If \cref{prob:entropy-inverse-q} holds for the primitive first-match residual families, then primitive asymptotic Q holds.  Consequently the Q cell is absorbed in the entropy--subfield frontier argument at every reserve larger than the logarithmic transition scale.
\end{corollary}

\begin{proof}
If primitive asymptotic Q failed, then by \cref{thm:logmoment-equivalence} there would be a logarithmic moment order \(r\) with \(\Gamma_r^{\rm prim}\ge\exp(\eta nr)\) along a subsequence for some \(\eta>0\).  Applying \cref{prob:entropy-inverse-q} gives either a paid cell, contrary to first-match removal, or a positive-density primitive subset in affine rank \(o(n)\).  The latter is impossible by \cref{prop:vandermonde-kills-low-rank} in the moment-curve range \(R=\Theta(n)\).  Thus \(M_{\rm prim}\le\exp(o(n))\), which is primitive asymptotic Q.
\end{proof}

\begin{definition}[Fourier-flat prefix leaf]\label{def:fourier-flat-prefix-leaf}
Let \(E\) be a finite field of cardinality \(Q\) and characteristic \(p\), let \(T\subseteq E\) have size \(N\), and let \(\Omega_m=\{M\subseteq T:|M|=m\}\).  For \(w<p\), define
\[
        \Psi(M)=\left(\sum_{y\in M}y,\sum_{y\in M}y^2,\ldots,\sum_{y\in M}y^w\right)\in E^w .
\]
Fix a nontrivial additive character \(\psi:E\to\mathbb C^\times\).  A sequence of such leaves is Fourier-flat if there is an integer \(L=L_N\) such that, for every nonzero \(\alpha=(\alpha_1,\ldots,\alpha_w)\in E^w\) and every \(1\le j\le m\) with \(p\nmid j\),
\[
        \left|\sum_{y\in T}\psi\left(j\sum_{i=1}^w\alpha_i y^i\right)\right|\le L,
\]
and, with
\[
        \theta=L+\left\lceil\frac{N-L}{p}\right\rceil ,
\]
one has
\[
        Q^w\binom{m+\theta-1}{m}\binom Nm^{-1}\le \exp(o(N)).
\]
It is strongly Fourier-flat if the last expression is \(o(1)\).
\end{definition}

\begin{theorem}[Fourier-flat leaves satisfy asymptotic Q]\label{thm:fourier-flat-q}
In a Fourier-flat prefix leaf,
\[
        \max_{\xi\in E^w}
        \left|\{M\in\Omega_m:\Psi(M)=\xi\}\right|
        \le
        \exp(o(N))\,Q^{-w}\binom Nm .
\]
If the leaf is strongly Fourier-flat, the normalized maximum is \(1+o(1)\).  Since \(w<p\), Newton identities identify the fibers of \(\Psi\) with the fibers of the first \(w\) locator coefficients.  Thus Fourier-flatness is a sufficient condition for the asymptotic Q input on that leaf.
\end{theorem}

\begin{proof}
For \(\xi\in E^w\), character orthogonality gives
\[
        |\Psi^{-1}(\xi)|
        =
        Q^{-w}\sum_{\alpha\in E^w}\psi(-\alpha\cdot\xi)
        \sum_{M\in\Omega_m}
        \psi\left(\sum_{y\in M}\sum_{i=1}^w\alpha_i y^i\right).
\]
The term \(\alpha=0\) is \(Q^{-w}\binom Nm\).  For \(\alpha\ne0\), apply the distinct-coordinate inclusion-exclusion sieve of Li--Wan \cite{LiWanSieve}.  A cycle of length divisible by \(p\) contributes at most \(N\), and every other cycle contributes at most \(L\) by hypothesis.  Hence the corresponding elementary-symmetric character sum is bounded by the coefficient of \(z^m\) in
\[
        (1-z)^{-L}(1-z^p)^{-(N-L)/p}.
\]
Replacing \((N-L)/p\) by \(\lceil(N-L)/p\rceil\) and using coefficientwise monotonicity gives the upper bound \(\binom{m+\theta-1}{m}\).  Therefore
\[
        |\Psi^{-1}(\xi)|
        \le
        Q^{-w}\binom Nm+\binom{m+\theta-1}{m}.
\]
The Fourier-flat inequality is exactly the assertion that the second term is at most \(\exp(o(N))\) times the random-model mean \(Q^{-w}\binom Nm\); the strongly Fourier-flat version gives \(1+o(1)\).  Finally, because \(w<p\), the first \(w\) elementary-symmetric locator coefficients and the first \(w\) power sums are related by invertible triangular Newton identities.
\end{proof}

\begin{corollary}[large-characteristic Fourier examples]\label{cor:large-characteristic-fourier-examples}
If \(T=E=\F_p\), \(m/p\to\rho\in(0,1)\), and \(w=o(\sqrt p)\), then the leaf is strongly Fourier-flat.  More generally the leaf is Fourier-flat whenever Weil's bound gives \(L=O(w\sqrt p)\) and
\[
        w\log p+\log\binom{m+O(w\sqrt p+p^{-1}N)-1}{m}
        \le
        \log\binom Nm+o(N).
\]
\end{corollary}

\begin{proof}
For \(T=\F_p\), every nonzero phase polynomial \(\sum_{i=1}^w\alpha_iX^i\) has degree \(<p\), so Weil's bound gives \(L\le(w-1)\sqrt p\) for all nontrivial modes.  If \(w=o(\sqrt p)\), then \(\theta=o(p)\), \(w\log p=o(p)\), and \(\log\binom{m+\theta-1}{m}=o(p)\), while \(\log\binom pm=\Theta(p)\) for fixed density.  The stated general criterion is the same calculation with the displayed bound for \(L\).
\end{proof}

\begin{remark}[proof skeleton for \cref{prob:entropy-inverse-q}]\label{rem:entropy-inverse-skeleton}
The expected proof has six steps.  First, a large \(\Gamma_r^{\rm prim}\) dyadically extracts popular fibers with normalized size \(\exp(\Omega(n))\) above random.  Second, comparing supports inside a popular fiber to a base support produces signed trades satisfying the first \(w\) moment equations.  Third, Tao's cyclic-prime uncertainty principle and the BCH/Vandermonde barrier \cite{Tao05} remove low-support primitive trades.  Fourth, an entropy Balog--Szemeredi--Gowers step should convert the logarithmic collision excess into an entropy-small-doubling population of trades.  Fifth, Green--Ruzsa/PFR-type structuralization should produce a low-rank coset progression.  Sixth, a slice-derivative lemma must push this structure back from trade sums to many individual moment-curve columns, after which \cref{prop:vandermonde-kills-low-rank} forces the excess into a paid cell.  The fifth and sixth steps are the nonstandard parts; this note records them as the remaining theorem, not as a completed proof.
\end{remark}

\begin{remark}[mass-aware collision bookkeeping]\label{rem:mass-aware-logmoment}
After first-match deletion, the primitive residual family may have total mass \(\tau<1\) relative to all \(m\)-subsets.  Consequently a proof of \cref{prob:entropy-inverse-q} may not import a full-mass lower bound such as \(\Gamma_r^{\rm prim}\ge1\) unless \(\tau=1\) is separately proved.  The finite proof route should therefore use an off-diagonal or falling-factorial collision moment, prove a dense-heavy-fiber hypothesis before invoking ordinary moments, or name sparse-heavy primitive fibers as an explicit residual cell.  This caveat does not affect \cref{thm:logmoment-equivalence}, whose proof uses only the displayed definitions and the inequality \(\sum_sN^{\rm prim}(s)\le\binom nm\).
\end{remark}

\begin{remark}[why finite adjacent Q needs high order]\label{rem:finite-moment-order}
At the deployed frontier depths, \(w\log_2|\B|\) is on the order of \(10^5\) to \(10^6\) bits, while the adjacent margins are \(22.2,22.0,3.3,3.1\) bits.  Thus a finite adjacent proof using only \cref{thm:moment-q} needs either very high moment order \(r\), an exchange-compression theorem after quotient routing, or a direct max-fiber certificate.  Fixed second or third moments can be asymptotically useful evidence, but they cannot by themselves fit the printed adjacent margins.  The raw mode-at-null shortcut is false by \cref{prop:mode-null-false}.
\end{remark}

\begin{proposition}[twist-orbit constancy]\label{prop:twist-orbit}
Assume \(D\) is stable under multiplication by a cyclic group \(H\) of order \(n\).  For every \(\zeta\in H\), the map
\[
        (z_1,\ldots,z_w)\longmapsto(\zeta z_1,\zeta^2z_2,\ldots,\zeta^wz_w)
\]
preserves prefix-fiber sizes.
\end{proposition}

\begin{proof}
The map \(M\mapsto \zeta M\) is a bijection on \(m\)-subsets of \(D\).  If
\[
        \ell_M(X)=X^m+\sum_{j=1}^m \lambda_j(M)X^{m-j},
\]
then \(\lambda_j(\zeta M)=\zeta^j\lambda_j(M)\).  Hence multiplication by \(\zeta\) sends the fiber over \(z\) bijectively to the fiber over \((\zeta^j z_j)_{j\le w}\).
\end{proof}

\begin{proposition}[mode-at-null is not the finite Q target]\label{prop:mode-null-false}
Let \(D=\F_{17}^{\times}\), \(m=9\), and \(w=1\).  For \(s\in\F_{17}\), set
\[
        N_9(s)=\#\{M\subseteq D:\ |M|=9,\ \sum_{x\in M}x=s\}.
\]
Then
\[
        N_9(0)=672,
        \qquad
        N_9(s)=673\quad(s\ne0).
\]
Thus the largest first-prefix fiber is not the null fiber.  The nonzero fibers form one primitive twist orbit under \(\F_{17}^{\times}\).
\end{proposition}

\begin{proof}
Let \(\omega=\exp(2\pi i/17)\).  Fourier inversion gives
\[
        N_9(s)=\frac{1}{17}\sum_{\lambda\in\F_{17}}\omega^{-\lambda s}
        [T^9]\prod_{x\in\F_{17}^{\times}}(1+T\omega^{\lambda x}).
\]
For \(\lambda=0\), the coefficient is \(\binom{16}{9}=11440\).  For \(\lambda\ne0\), multiplication by \(\lambda\) permutes \(D\), and
\[
        \prod_{x\in D}(1+T\omega^{\lambda x})
        =
        \prod_{y=1}^{16}(1+T\omega^y)
        =\frac{1+T^{17}}{1+T}
        =1-T+T^2-\cdots+T^{16}.
\]
The coefficient of \(T^9\) in this last polynomial is \(-1\).  Therefore
\[
        N_9(0)=\frac{11440+16(-1)}{17}=672,
        \qquad
        N_9(s)=\frac{11440-(-1)}{17}=673
        \quad(s\ne0),
\]
because \(\sum_{\lambda\ne0}\omega^{-\lambda s}=-1\) for \(s\ne0\).  Since \(w=1\), any nonzero target has twist stabilizer \(\gcd(16,1)=1\), so the nonzero targets form a full primitive orbit by \cref{prop:twist-orbit}.
\end{proof}

\begin{proposition}[twist-orbit moment amplification]\label{prop:q-orbit-moment}
Assume \(D\) is stable under multiplication by a cyclic group \(H\) of order \(n\).  For \(z=(z_1,\ldots,z_w)\in\B^w\), put
\[
        A(z)=\{1\le j\le w:z_j\ne0\},
        \qquad
        s(z)=
        \begin{cases}
        n, & A(z)=\varnothing,\\
        \gcd(n,A(z)), & A(z)\ne\varnothing .
        \end{cases}
\]
Let
\[
        R(z)=|\B|^w \frac{|\Fib_w(z)|}{\binom nm},
        \qquad
        \Gamma_r=|\B|^{w(r-1)}
        \sum_y\left(\frac{|\Fib_w(y)|}{\binom nm}\right)^r
        \quad(r\ge2).
\]
Then
\[
        \Gamma_r\ge \frac{n}{s(z)}\,\frac{R(z)^r}{|\B|^w}.
\]
Consequently, if \(\Gamma_r\le G\), then
\[
        R(z)\le\left(\frac{G|\B|^w s(z)}{n}\right)^{1/r}.
\]
In particular, a primitive prefix value with \(\gcd(n,A(z))=1\) gains a full factor \(n\) in the moment test.
\end{proposition}

\begin{proof}
By \cref{prop:twist-orbit}, every prefix value in the orbit
\[
        (\zeta z_1,\zeta^2z_2,\ldots,\zeta^wz_w),
        \qquad \zeta\in H,
\]
has the same fiber size as \(z\).  The stabilizer consists exactly of those \(\zeta\in H\) satisfying \(\zeta^j=1\) for every \(j\in A(z)\), so its size is \(s(z)\).  Hence the orbit has size \(n/s(z)\).  The contribution of this orbit alone to \(\Gamma_r\) is
\[
        |\B|^{w(r-1)}\frac{n}{s(z)}
        \left(\frac{|\Fib_w(z)|}{\binom nm}\right)^r
        =
        \frac{n}{s(z)}\frac{R(z)^r}{|\B|^w}.
\]
This gives the first inequality, and the second follows by rearranging.
\end{proof}

\begin{proposition}[composite prefix directions descend]\label{prop:composite-descend}
Let \(S\) be a multiplicative coset of order \(N\) in \(\F_p^\times\), let \(e\ge1\), put \(c=\gcd(e,N)\), and let \(\psi:\F_p\to\mathbb C^\times\) be a nontrivial additive character.  Suppose \(g(x)=h(x^e)\).  If \(S_e=\{a^e:a\in S\}\), then for every \(m\)
\[
        \sum_{\substack{A\subseteq S\\ |A|=m}}
        \psi\!\left(\sum_{a\in A}g(a)\right)
        =
        [T^m]\prod_{b\in S_e}\bigl(1+T\psi(h(b))\bigr)^c .
\]
Thus a composite Fourier direction contributes through the image coset \(S_e\)
with power-map multiplicity \(c\); when \(c>1\) it is a quotient-scale object
rather than primitive mass.
\end{proposition}

\begin{proof}
The left side is the coefficient of \(T^m\) in
\[
        \prod_{a\in S}(1+T\psi(g(a))).
\]
Since \(g(a)=h(a^e)\) is constant on fibers of \(a\mapsto a^e\), and the power map has exactly \(c=\gcd(e,N)\) preimages over each point of its image on the coset \(S\), the product factors as displayed.
\end{proof}

\begin{theorem}[proved head-depth prefix flatness]\label{thm:head-flatness}
Let \(p\) be prime, \(H\le\F_p^\times\) have order \(N\), put \(d=(p-1)/N\), and let \(S=\alpha H\).  Let \(P\subseteq S\) have size \(\ell\), let \(1\le w<N-\ell\) satisfy \(wd<p\), and put \(m'=m-\ell\).  For \(\vec z\in\F_p^w\), define
\[
        \mathcal N_w^P(m,\vec z)=
        \#\{A:P\subseteq A\subseteq S,\ |A|=m,\ (e_1(A),\ldots,e_w(A))=\vec z\}.
\]
Then
\[
        \left|\mathcal N_w^P(m,\vec z)-p^{-w}\binom{N-\ell}{m'}\right|
        \le
        p^{w/2}\binom{w\lceil\sqrt p\rceil+\ell+m'-1}{m'} .
\]
For \(w=1\) the prefactor \(p^{1/2}\) may be omitted.
\end{theorem}

\begin{proof}
Write \(A=P\cup A'\), \(A'\subseteq S\setminus P\).  Since \(\ell_A=\ell_P\ell_{A'}\), the first \(w\) elementary symmetric functions of \(A\) are an affine unipotent triangular transform of those of \(A'\).  Thus planted fibers have the same sizes as unplanted fibers over transformed values, except that the ambient set has size \(N-\ell\).

By Fourier inversion over \(\F_p^w\), the deviation from the mean is controlled by
\[
        E(\vec t)=
        \sum_{\substack{A'\subseteq S\setminus P\\ |A'|=m'}}
        \exp\left(\frac{2\pi i}{p}\sum_{j=1}^w t_j e_j(A')\right),
        \qquad \vec t\ne0.
\]
Let \(j_0\) be the largest index with \(t_{j_0}\ne0\).  Newton identities rewrite \(\sum_jt_je_j(A')\) as a polynomial in the power sums \(p_1(A'),\ldots,p_{j_0}(A')\) that is nontrivially linear in the last variable.  A second Fourier inversion therefore reduces \(E(\vec t)\) to elementary symmetric functions of the numbers
\[
        \exp\left(\frac{2\pi i}{p}g(a)\right),
        \qquad a\in S\setminus P,
\]
where \(g\in\F_p[X]\) is nonconstant of degree at most \(w\).

For every nonzero integer \(r\) modulo \(p\), the power sum of this multiset is bounded by \(w\lceil\sqrt p\rceil+\ell\).  Indeed, parameterize the coset as \(S=\{\alpha y^d:y\in\F_p^\times\}\).  Since \(\deg_y(g(\alpha y^d))\le wd<p\), Weil's bound for nonconstant additive character sums \cite{Weil48} gives \(w\sqrt p\) up to harmless endpoint terms; removing the \(\ell\) planted points costs at most \(\ell\).

Newton's recursion for elementary symmetric functions is coefficientwise dominated by the recursion for \((1-u)^{-(w\lceil\sqrt p\rceil+\ell)}\).  Hence every reduced Fourier coefficient is bounded by
\[
        \binom{w\lceil\sqrt p\rceil+\ell+m'-1}{m'}.
\]
The second Fourier transform contributes at most \(p^{j_0/2}\) by Parseval, and \(j_0\le w\).  Summing the first Fourier inversion and dividing by \(p^w\) gives the stated \(p^{w/2}\) bound.  When \(w=1\), the Newton-linearization step is absent and the Parseval factor is unnecessary.
\end{proof}

\begin{corollary}[the solved head of Q]\label{cor:head-q}
For the KoalaBear identity-scale subgroup \(N=2^{21}\), \(d=1016\), and \(\ell=0\), \cref{thm:head-flatness} proves two-sided prefix flatness at the deployed subset sizes \(m=1116047\) and \(m=1116046\) for exactly \(w\le21\).  At the head rungs \(m=K+w\), it is nonvacuous for exactly \(w\le22\).  In these ranges every prefix fiber equals
\[
        \binom Nm p^{-w}(1+\varepsilon)
\]
with \(|\varepsilon|\) exponentially small in the recorded margin.
\end{corollary}

\begin{proof}
Substitute the displayed row parameters into \cref{thm:head-flatness}.  The nonvacuous range is exactly the range in which the error term is smaller than the main term.  Using the standard entropy sandwich
\[
        \frac{2^{M H_2(r/M)}}{M+1}\le \binom Mr\le 2^{M H_2(r/M)}
\]
turns this into explicit finite inequalities in \(p,N,m,w\), which are the head-depth inequalities recorded in the v13 raw manuscript.
\end{proof}

\begin{remark}[why head-depth Q does not close the dense frontier]\label{rem:head-does-not-close-q}
\Cref{thm:head-flatness} pays one square-root loss per active power-sum direction.  This is why it reaches only \(w=21\) or \(22\) in the KoalaBear head regimes, while the deployed frontier rows have \(w\approx 67466\).  The next Q target is therefore not a uniform per-character bound.  Composite Fourier directions with nontrivial power-map multiplicity are quotient-scale objects by \cref{prop:composite-descend}; the \(c=1\) composite directions and primitive character maxima still face the usual Parseval square-root barrier.  The remaining statement must bound the maximum prefix fiber itself, after quotient and planted strata have been paid.
\end{remark}

\begin{proposition}[Fourier audit for Q]\label{prop:fourier-audit}
Let \(D\subseteq\F_p\), let \(\Phi_w:\binom{D}{m}\to\F_p^w\) be the prefix map, and write \(N(z)=|\Phi_w^{-1}(z)|\) and \(\overline F=\binom{|D|}{m}p^{-w}\).  For \(t\in\F_p^w\), define
\[
        E(t)=
        \sum_{\substack{M\subseteq D\\ |M|=m}}
        \exp\left(\frac{2\pi i}{p}\,t\cdot\Phi_w(M)\right).
\]
Then
\[
        N(z)-\overline F
        =
        p^{-w}\sum_{t\ne0}
        \exp\left(-\frac{2\pi i}{p}\,t\cdot z\right)E(t),
        \qquad
        \sum_{t\ne0}|E(t)|^2=p^w\sum_z |N(z)-\overline F|^2 .
\]
Consequently a proof of Q that uses only the triangle inequality and a uniform nonzero-character estimate must reach the scale
\[
        \max_{t\ne0}|E(t)|\le \poly(n)\binom{|D|}{m}p^{-w}.
\]
Uniform square-root estimates of order \(\binom{|D|}{m}p^{-w/2}\) are short by a factor \(p^{w/2}\) and cannot close the frontier-depth Q input without additional joint cancellation or max-fiber structure.
\end{proposition}

\begin{proof}
The function \(N\) on the finite abelian group \(\F_p^w\) has Fourier transform \(E\).  Fourier inversion gives the first displayed identity, and the zero-frequency term is exactly \(\overline F\).  Plancherel applied to \(N-\overline F\) gives the second identity, since the zero Fourier coefficient of \(N-\overline F\) is zero.

If the first formula is bounded by triangle inequality, then
\[
        |N(z)-\overline F|\le p^{-w}(p^w-1)\max_{t\ne0}|E(t)|<\max_{t\ne0}|E(t)|.
\]
Q asks, after quotient and planted strata have been removed, for \(N(z)\le\poly(n)\overline F\).  Since \(\overline F=\binom{|D|}{m}p^{-w}\), the required individual-character scale is the one displayed.  The square-root scale is larger by \(p^{w/2}\).
\end{proof}

\begin{proposition}[quotient pullback for shift pairs]\label{prop:sp-pullback}
Assume \(D\) is a multiplicative coset and \(c\mid n\).  Suppose \(A(X)=A_0(X^c)\) and \(B(X)=B_0(X^c)\) are a depth-\(w\) shift pair of degree \(e=ce_0\), meaning \(\deg(A-B)\le e-w-1\).  Then \((A_0,B_0)\) is a shift pair on the quotient coset at depth
\[
        w_c=\left\lceil\frac{w+1}{c}\right\rceil-1.
\]
Conversely, any quotient pair at depth \(w_c\) pulls back to a depth-\(w\) pair upstairs.
\end{proposition}

\begin{proof}
The inequality \(c\deg(A_0-B_0)\le ce_0-w-1\) is equivalent to
\[
        \deg(A_0-B_0)\le e_0-\left\lceil\frac{w+1}{c}\right\rceil,
\]
which is precisely the depth-\(w_c\) condition.  The converse is the same integer inequality read in the other direction.
\end{proof}

\begin{definition}[coefficient scale]\label{def:coefficient-scale}
Assume \(D=\alpha H\subseteq\B^\times\), where \(H\) is cyclic of order \(n\) and \(\operatorname{char}\B\nmid n\).  For a monic degree-\(e\) polynomial
\[
        L(X)=X^e+\sum_{j=1}^e \lambda_jX^{e-j}
\]
split over \(D\), define
\[
        s(L)=\gcd\bigl(n,e,\{j:\lambda_j\ne0\}\bigr).
\]
For a pair \((A,B)\) of monic split degree-\(e\) polynomials, put \(s(A,B)=\gcd(s(A),s(B))\).
\end{definition}

\begin{lemma}[coefficient scale detects quotient support]\label{lem:coeff-scale}
Let \(S\subseteq D\), \(|S|=e\), and \(L_S=\ell_S\).  For \(c\mid n\), the following are equivalent:
\[
        S\text{ is a union of fibers of }x\mapsto x^c,\qquad
        L_S(X)=L_c(X^c)\text{ for a monic }L_c,\qquad
        c\mid s(L_S).
\]
In this case \(L_c\) is uniquely determined and split over the quotient coset \(D_c=\{x^c:x\in D\}\).
\end{lemma}

\begin{proof}
If \(S\) is a union of fibers of \(x\mapsto x^c\), then for \(\overline S=\{x^c:x\in S\}\) one has
\[
        L_S(X)=\prod_{y\in\overline S}(X^c-y)=L_c(X^c).
\]
This proves the first implication and uniqueness.  The second displayed condition is equivalent to \(c\mid e\) and to all nonzero coefficient gaps \(j\) being divisible by \(c\), which is exactly \(c\mid s(L_S)\).  Conversely, if \(L_S=L_c(X^c)\), then every root \(x\in D\) brings with it the whole fiber of \(x\mapsto x^c\), since \(X^n-\alpha^n\) is squarefree and \(c\mid n\).  Thus \(S\) is a union of quotient fibers.
\end{proof}

\begin{theorem}[maximal quotient extraction for shift pairs]\label{thm:coeff-quotient-extract}
Let \((A,B)\) be a depth-\(w\) shift pair of degree \(e\) over \(D\):
\[
        A=\ell_S,\qquad B=\ell_T,\qquad S\cap T=\varnothing,\qquad
        \deg(A-B)\le e-w-1.
\]
Let \(c=s(A,B)\).  If \(c>1\), then \((A,B)\) is quotient-borne.  More precisely, there are unique monic locators \(A_c,B_c\), split over \(D_c\), such that
\[
        A(X)=A_c(X^c),\qquad B(X)=B_c(X^c),
\]
and \((A_c,B_c)\) is a depth
\[
        w_c=\left\lceil\frac{w+1}{c}\right\rceil-1
\]
shift pair on \(D_c\).  If \(c\) is maximal, then the descended pair has common coefficient scale \(1\) relative to the quotient length \(n/c\).
\end{theorem}

\begin{proof}
\Cref{lem:coeff-scale} gives the unique quotient locators.  The depth statement is exactly \cref{prop:sp-pullback}.  For maximality, suppose the descended pair had a nontrivial common coefficient scale \(d>1\) relative to \(n/c\).  Then \(A_c=A_{cd}(Y^d)\) and \(B_c=B_{cd}(Y^d)\), so \(A=A_{cd}(X^{cd})\) and \(B=B_{cd}(X^{cd})\), contradicting maximality of \(c\).
\end{proof}

\begin{corollary}[primitive exclusion by coefficient support]\label{cor:primitive-coeff-exclusion}
Every primitive shift pair has
\[
        \gcd\bigl(n,e,\{j:\lambda_j(A)\ne0\},\{j:\lambda_j(B)\ne0\}\bigr)=1.
\]
In particular, the top-stratum binomial prototypes \(A=X^{w+1}-\beta\), \(B=X^{w+1}-\beta'\), when split over \(D\), are quotient-borne and belong to the quotient ledger rather than the primitive SP cell.
\end{corollary}

\begin{proof}
The displayed gcd is \(s(A,B)\).  If it is \(>1\), \cref{thm:coeff-quotient-extract} extracts a quotient pair, so the pair is not primitive.  For the binomial prototypes, the only nonzero coefficient gap is \(w+1=e\), and splitting over the order-\(n\) coset forces \(w+1\mid n\); hence the common coefficient scale is \(w+1\).
\end{proof}

\begin{proposition}[top-stratum quotient sieve]\label{prop:top-stratum-quotient-sieve}
Assume \(D=\alpha H\subseteq\B^\times\), where \(H\) is cyclic of order \(n\) and \(\operatorname{char}\B\nmid n\).  Let \((A,B)=(\ell_S,\ell_T)\) be a depth-\(w\) shift pair of degree \(e=w+1\) over \(D\), with \(S\cap T=\varnothing\), and put \(c=s(A,B)\).  If \(c>1\), then \(c\mid w+1\), and maximal quotient extraction gives unique monic locators \(A_c,B_c\), split over \(D_c=\{x^c:x\in D\}\), such that
\[
        A(X)=A_c(X^c),\qquad B(X)=B_c(X^c),
\]
with
\[
        \deg A_c=\deg B_c=\frac{w+1}{c},
        \qquad
        A_c-B_c\in\B^\times .
\]
Thus every non-primitive top-stratum shift pair is paid by a top-stratum constant-shift cell on a quotient coset.  The primitive top-stratum residual consists only of constant-shift pairs \(A,A-\lambda\), \(\lambda\in\B^\times\), with \(s(A,A-\lambda)=1\).
\end{proposition}

\begin{proof}
Since \(c=s(A,B)\), the definition of coefficient scale gives \(c\mid e=w+1\).  By \cref{thm:coeff-quotient-extract}, there are unique quotient locators \(A_c,B_c\) split over \(D_c\), and the descended pair has depth
\[
        w_c=\left\lceil \frac{w+1}{c}\right\rceil-1
        =\frac{w+1}{c}-1 .
\]
Its degree is \(e_c=e/c=(w+1)/c\).  Therefore
\[
        \deg(A_c-B_c)\le e_c-w_c-1=0 .
\]
The constant is nonzero: if \(A_c=B_c\), then \(A=B\), contradicting disjointness of the positive-degree root sets.  Hence \(A_c-B_c\in\B^\times\).

At the upstairs top stratum \(e=w+1\), the shift-pair condition itself says \(\deg(A-B)\le0\), so every top-stratum pair is a constant-shift pair.  The quotient-borne ones are exactly those with \(s(A,B)>1\), by \cref{thm:coeff-quotient-extract}.  The remaining top-stratum cell is therefore the coefficient-primitive constant-shift cell \(s(A,B)=1\).
\end{proof}

For later reference, split the off-diagonal sum in \cref{prop:second-moment} into quotient-pulled-back pairs and primitive pairs by assigning every periodic pair to its maximal common quotient scale.  Put
\[
        \overline F=\binom{n}{m}|\B|^{-w},
        \qquad
        \Gamma_2=|\B|^w\sum_z\left(\frac{|\Fib_w(z)|}{\binom{n}{m}}\right)^2 .
\]

\begin{proposition}[exact primitive \(L^2\) ledger]\label{prop:gamma2-ledger}
With the first-match quotient assignment above,
\[
        \Gamma_2
        =
        \frac{1}{\overline F}
        +\frac{\sum_e P_e^{\rm quot}}{\binom{n}{m}\overline F}
        +\frac{\sum_e P_e^{\rm prim}}{\binom{n}{m}\overline F},
\]
where \(P_e^{\rm quot}\) and \(P_e^{\rm prim}\) are the quotient and primitive parts of the \(e\)-th off-diagonal term in \cref{prop:second-moment}.  Thus SP is exactly the task of bounding the last displayed summand after quotient-pulled-back pairs have been paid.
\end{proposition}

\begin{proof}
Divide the identity in \cref{prop:second-moment} by \(\binom{n}{m}^2|\B|^{-w}\).  The diagonal contributes \(1/\overline F\); the off-diagonal terms split by the first-match quotient/primitive assignment.
\end{proof}

\begin{proposition}[slope elimination]\label{prop:slope-elimination}
Fix \(K\) and \(m\ge K\).  Let \(u+zv\) be an affine line of received words and let \(T\subseteq D\) have \(|T|=m\).  Put \(w=m-K\) and
\[
        S_r(h,T)=\sum_{x\in T}\frac{h(x)x^r}{\ell_T'(x)},\qquad 0\le r<w.
\]
Then \(u+zv\) restricts on \(T\) to a polynomial of degree \(<K\) if and only if
\[
        (S_r(u,T))_{r<w}+z(S_r(v,T))_{r<w}=0.
\]
If \(T\) is not a common support for both \(u\) and \(v\), then it supports at most one finite slope \(z\).
\end{proposition}

\begin{proof}
The Lagrange interpolant of a word \(h:T\to\F\) has degree \(<K\) if and only if its top \(w\) coefficients vanish; the displayed \(S_r\)'s are the usual residue formula for those top coefficients, up to an invertible triangular change of basis.  Applying this to \(h=u+zv\) gives the vector equation.  If the second vector is zero, the support is either common or carries no finite slope.  If it is nonzero, any solution \(z\) is forced by a nonzero coordinate, so there is at most one.
\end{proof}

\begin{proposition}[interpolation-lattice split-pencil reduction]\label{prop:lattice-split}
Fix \(K\) and an agreement \(m\ge K\), and write \(|D|=n\).  For a word \(U:D\to\F\), define
\begin{align*}
        M_U&=\{(W,N)\in\F[X]^2:
        W(x)U(x)=N(x)\text{ for all }x\in D\},\\
        \wdeg(W,N)&=\max(\deg W,\deg N-(K-1)).
\end{align*}
Let \(\omega=n-m\).  Agreement supports of size \(m\) for \(C=\RS[\F,D,K]\) are in bijection with pairs \((W,N)\in M_U\) such that \(W\) is monic, \(D\)-split, \(\deg W=\omega\), \(\wdeg(W,N)\le\omega\), and \(W\mid N\).  If a shifted weak-Popov basis \(g_i=(W_i,N_i)\) of \(M_U\) has profile \((d_1,d_2)\), every such pair has the form \(Ag_1+Bg_2\) with \(\deg A\le\omega-d_1\) and \(\deg B\le\omega-d_2\).  When \(W=AW_1+BW_2\) is \(D\)-split, the divisibility \(W\mid N\) is automatic.
\end{proposition}

\begin{proof}
If \(U\) agrees with a degree-\(<K\) polynomial \(c\) on \(T\), take \(W=\ell_{D\setminus T}\) and \(N=Wc\).  Conversely, a pair with \(W\mid N\) gives \(N=Wc\), the weighted-degree bound gives \(\deg c<K\), and the relation \(W(U-c)=0\) on \(D\) gives agreement on \(D\setminus Z(W)\).  The weak-Popov representation is the predictable-degree property of a shifted basis of the rank-two interpolation module.  Finally, if \(W\) is \(D\)-split, then at every root \(x\in D\) of \(W\) the module relation gives \(N(x)=W(x)U(x)=0\); since \(W\) is squarefree, \(W\mid N\).
\end{proof}

\begin{theorem}[near-rational lattice dichotomy]\label{thm:near-rational}
Let \(C=\RS[\F,D,K]\), let \(U:D\to\F\), let \(m\ge K\), put \(w=m-K\) and \(\omega=n-m\), and assume \(2w\le n-K\).  Let \(d_1=d_1(U)\) be the smallest shifted degree of a nonzero element of \(M_U\), and let
\[
        \Cen(U;m)=\#\{T\subseteq D:\ |T|=m,\ U|_T\text{ is explained by a degree-}<K\text{ polynomial}\}.
\]
Then exactly one of the following alternatives holds.
\begin{enumerate}[label=\textup{(\roman*)}]
\item If \(d_1\le w\), then either \(U\) is within Hamming distance exactly \(d_1\) of a unique codeword \(c\in C\), in which case
\[
        \Cen(U;m)=\binom{n-d_1}{m},
\]
or \(\Cen(U;m)=0\).
\item If \(d_1\ge w+1\), every census element is represented in the two-generator family
\[
        Ag_1+Bg_2,\qquad
        \deg A\le\omega-d_1,\quad \deg B\le\omega-d_2,
\]
where \(g_1,g_2\) is a shifted weak-Popov basis of \(M_U\), \(d_2=n-K+1-d_1\).  The parameter space has dimension \(\omega-w+1\).
\end{enumerate}
\end{theorem}

\begin{proof}
Let \(g_1=(W_1,N_1)\), \(g_2=(W_2,N_2)\) be a shifted weak-Popov basis with row degrees \(d_1\le d_2\).  The determinant degree gives \(d_1+d_2=n-K+1\).  If \(d_1\le w\), then \(d_2\ge n-K+1-w=\omega+1\).  By predictable degrees, any element of shifted degree at most \(\omega\) must be a multiple \(Ag_1\).  Therefore the census is empty unless \(W_1\), after scaling, is a monic squarefree divisor of \(\ell_D\), \(W_1\mid N_1\), and \(c=N_1/W_1\) has degree \(<K\).  In the nonempty case, \(c\) agrees with \(U\) outside the roots of \(W_1\).  Minimality of \(d_1\) forces the disagreement set to be exactly those roots; otherwise its error-locator pair would be a nonzero lattice element of smaller shifted degree.  The codeword is unique because two codewords within distance \(w\) of \(U\) would differ in at most \(2w\le n-K\) positions, contradicting the Reed--Solomon distance \(n-K+1\).  The valid supports are precisely the \(m\)-subsets of the \(n-d_1\) agreement positions, giving the binomial count.  If any required split/divisibility/codeword condition fails, the census is empty.

If \(d_1\ge w+1\), both weak-Popov rows may occur.  Predictable degrees give the multiplier bounds \(\deg A\le\omega-d_1\) and \(\deg B\le\omega-d_2\).  The number of coefficients in \(A\) and \(B\) is
\[
        (\omega-d_1+1)+(\omega-d_2+1)
        =2\omega+2-(n-K+1)=\omega-w+1.
\]
\end{proof}

\begin{corollary}[near-rational slopes are paid outside the balanced core]\label{cor:near-rational-line}
For \(C=\RS[\F,D,K]\) and an affine line \(u+zv\), suppose two distinct slopes \(z_1,z_2\) have \(d_1(u+z_iv)\le w\) and nonzero size-\(m\) census.  Then \(v\) is within distance \(2w\) of a codeword and \(u\) is within distance \(3w\) of a codeword.  Thus this line belongs to a common-proximity paid branch; after that branch is removed, at most one near-rational finite slope remains, and all other residual slopes satisfy \(d_1(u+zv)\ge w+1\).
\end{corollary}

\begin{proof}
By \cref{thm:near-rational}, \(u+z_i v=c_i+\eta_i\), with \(c_i\in C\) and \(\operatorname{wt}(\eta_i)\le w\), for \(i=1,2\).  Since \(z_1\ne z_2\),
\[
        v=\frac{c_1-c_2}{z_1-z_2}+\frac{\eta_1-\eta_2}{z_1-z_2}
\]
is within distance at most \(2w\) of a codeword.  Substituting into \(u=(u+z_1v)-z_1v\), the word \(u\) is within distance at most \(w+2w=3w\) of a codeword.  Such a line is therefore not part of the primitive balanced-core residual once common-proximity/tangent cells are paid.  If two near-rational slopes survived there, the preceding argument would put the line in the paid branch; hence at most one can remain.
\end{proof}

\begin{theorem}[deficiency-one split-test eliminant]\label{thm:deficiency-one-eliminant}
Let \(\Lambda_D\in\F[X]\) be squarefree of degree \(n\).  Let \(M(Z)=M_0+ZM_1\) be a \(j\times(j+1)\) matrix pencil over \(\F\).  For \(0\le i\le j\), let \(c_i(Z)\) be the signed maximal minor obtained by deleting column \(i\), and set
\[
        L_Z(X)=\sum_{i=0}^j c_i(Z)X^i .
\]
Assume not all maximal minors vanish identically.  If \(c_j\not\equiv0\), perform pseudo-division over \(\F[Z]\):
\[
        c_j(Z)^{n-j+1}\Lambda_D(X)=Q(Z,X)L_Z(X)+R(Z,X),
        \qquad \deg_XR<j .
\]
If some coefficient \(R_m(Z)\) is not identically zero, then every parameter \(z\) for which either \(M(z)\) drops rank, or \(M(z)\) has full rank with \(c_j(z)=0\), or \(M(z)\) has full rank with \(c_j(z)\ne0\) and \(L_z\mid\Lambda_D\), is a root of the nonzero polynomial \(c_j(Z)R_m(Z)\).  Its degree is at most \(j+(n-j+1)j\).  If \(c_j\equiv0\), the top chart is empty.  If \(c_j\not\equiv0\) and all coefficients of \(R\) vanish identically, then every top-chart parameter passes the split test \(L_z\mid\Lambda_D\), so this chart is a named residual branch rather than a paid bounded chart.
\end{theorem}

\begin{proof}
Each maximal minor is a determinant of a \(j\times j\) matrix with affine-linear entries in \(Z\), so it has degree at most \(j\).  At a full-rank value \(z\), the cofactor vector \(c(z)=(c_0(z),\ldots,c_j(z))\) spans the one-dimensional kernel of \(M(z)\).  Thus \(c_j(z)=0\) is the full-rank chart where the candidate locator has degree \(<j\), while \(c_j(z)\ne0\) is the top chart.

Pseudo-division by the degree-\(j\) polynomial \(L_Z\) takes \(n-j+1\) elimination steps.  Each step multiplies the current remainder by \(c_j\) and subtracts a shifted multiple of \(L_Z\), so every coefficient of \(R\) has \(Z\)-degree at most \((n-j+1)j\).

Let \(z\) be a top-chart parameter with \(L_z\mid\Lambda_D\).  Since \(c_j(z)\ne0\), substituting \(Z=z\) into the pseudo-division identity gives that \(L_z\) divides \(R(z,X)\).  But \(\deg_XR(z,X)<j=\deg L_z\), hence \(R(z,X)=0\); in particular \(R_m(z)=0\).  If \(M(z)\) drops rank, all maximal minors vanish, including \(c_j(z)\).  If \(M(z)\) has full rank but \(c_j(z)=0\), again \(c_j(z)=0\).  These two classes are therefore roots of \(c_jR_m\).  The degree bound follows from the preceding paragraph and \(\deg c_j\le j\).  Finally, if all coefficients of \(R\) vanish identically, then the same pseudo-division identity gives \(L_z\mid\Lambda_D\) for every top-chart \(z\), and no eliminant is available.
\end{proof}

\begin{lemma}[split annihilator dictionary]\label{lem:split-annihilator}
Let \(D\subseteq\F\) have size \(n\), let \(\Lambda_D(X)=\prod_{x\in D}(X-x)\), and put \(v_x=\Lambda_D'(x)^{-1}\).  Let \(L\in\F[X]\) be split over \(D\), with root set \(T\), \(\gamma=|T|=\deg L\), and let \(\tau\ge1\) satisfy \(\tau+\gamma\le n\).  For a word \(y:D\to\F\), the conditions
\[
        \sum_{x\in D}v_x\,y(x)L(x)x^m=0
        \qquad(0\le m<\tau)
\]
hold if and only if
\[
        y\in \RS[\F,D,n-\tau-\gamma]+\F^T,
\]
where \(\F^T\) denotes words supported on \(T\).  The sum is direct.
\end{lemma}

\begin{proof}
First note the residue identity: for every polynomial \(h\) with \(\deg h\le n-2\),
\[
        \sum_{x\in D}v_xh(x)=0.
\]
Indeed, \(h/\Lambda_D\) has no residue at infinity, and its finite residues are \(v_xh(x)\).

If \(y\) is supported on \(T\), then \(y(x)L(x)=0\) for all \(x\in D\).  If \(y=p|_D\) with \(\deg p<n-\tau-\gamma\), then for \(m<\tau\) the polynomial \(pLX^m\) has degree at most \(n-2\), so the residue identity gives the displayed annihilator equations.  Thus the right side lies in the solution space.

Conversely, the linear map
\[
        y\longmapsto
        \left(\sum_{x\in D}v_x\,y(x)L(x)x^m\right)_{0\le m<\tau}
\]
has zero columns on \(T\).  On \(D\setminus T\), the nonzero scalars \(v_xL(x)\) multiply a \(\tau\times(n-\gamma)\) Vandermonde matrix on distinct points, hence the rank is \(\tau\).  The solution space therefore has dimension \(n-\tau\).  The right side has dimension \((n-\tau-\gamma)+\gamma=n-\tau\), because a degree-\(<n-\tau-\gamma\) polynomial vanishing on \(D\setminus T\) has at least \(n-\gamma\) roots, more than its degree unless it is zero.  The containment of equal-dimensional spaces proves the equivalence and directness.
\end{proof}

\begin{lemma}[constant divisor rigidity]\label{lem:constant-divisor}
Let \(P\in\F[X]\) be nonzero and let \(L\in\F(Z)[X]\) be monic.  If \(L\mid P\) in \(\F(Z)[X]\), then \(L\in\F[X]\) and \(L\mid P\) already in \(\F[X]\).
\end{lemma}

\begin{proof}
Over \(\overline{\F}(Z)\), all roots of \(P\) are constants.  Since \(L\) is a monic divisor of \(P\), it is a product of some of the linear factors \(X-\xi\) with \(\xi\in\overline{\F}\).  Thus the coefficients of \(L\) lie in \(\overline{\F}\).  They also lie in \(\F(Z)\), and the relative algebraic closure of \(\F\) in \(\F(Z)\) is \(\F\); hence the coefficients lie in \(\F\).  Division with remainder in \(\F[X]\) agrees with division in \(\F(Z)[X]\), so the remainder is zero.
\end{proof}

\begin{proposition}[split charts are tangent-borne]\label{prop:split-chart-tangent}
Let \(C=\RS[\F,D,k]\), set \(j=n-a\), and assume \(j=a-k\).  Let \(M(Z)\) be the deficiency-one syndrome pencil of a received line \(f+Zg\).  In the notation of \cref{thm:deficiency-one-eliminant}, assume \(c_j\not\equiv0\) and all pseudo-division remainders vanish identically.  Then there is a fixed \(j\)-set \(T^*\subseteq D\) such that \(f\) and \(g\) each agree with codewords of \(C\) on \(D\setminus T^*\).  Consequently, at every threshold \(a'\ge a\), the line has at most \(j\) MCA-bad finite slopes, all among the tangent ratios
\[
        \begin{aligned}
        &-\frac{e_f(x)}{e_g(x)},\\
        &x\in T^*,\quad e_g(x)\ne0,
        \end{aligned}
\]
where \(e_f,e_g\) are the two error words supported on \(T^*\).  It has no CA-bad slopes whenever the internal radius is at least \(j/n\).
\end{proposition}

\begin{proof}
The vanishing of all pseudo-division remainders says that the normalized top-chart locator
\[
        \widehat L_Z(X)=c_j(Z)^{-1}\sum_{i=0}^j c_i(Z)X^i
\]
divides \(\Lambda_D\) in \(\F(Z)[X]\).  By \cref{lem:constant-divisor}, \(\widehat L_Z=L^*\in\F[X]\) is a fixed monic divisor of \(\Lambda_D\).  Let \(T^*\) be its root set in \(D\), so \(|T^*|=j\).

Because the cofactor vector \(c(Z)\) is a kernel vector of the syndrome pencil, the identity \(M(Z)c(Z)\equiv0\) separates into the two identities saying that \(L^*\) annihilates the exact-agreement syndromes of \(f\) and of \(g\).  Applying \cref{lem:split-annihilator} with \(\tau=j\) and \(\gamma=j\) gives \(f=c_f+e_f\) and \(g=c_g+e_g\), with \(c_f,c_g\in\RS[\F,D,n-2j]=C\) and \(e_f,e_g\) supported on \(T^*\).  Here \(n-2j=k\) by \(j=n-a=a-k\).

Let \(z\) be MCA-bad at agreement \(a'\ge a\), with witness support \(S\) and explaining codeword \(c\) for \(f+zg\).  On \(S\setminus T^*\) the codewords \(c\) and \(c_f+zc_g\) agree.  This set has size at least \(a'-j\ge a-j=k\), so they are the same codeword.  Thus on \(S\cap T^*\) one has \(e_f+ze_g=0\).  If the pair were simultaneously explained on \(S\), it would be explained by \((c_f,c_g)\), so badness forces some \(x\in S\cap T^*\) with \((e_f(x),e_g(x))\ne(0,0)\).  At such an \(x\), the equality \(e_f(x)+ze_g(x)=0\) forces \(e_g(x)\ne0\) and \(z=-e_f(x)/e_g(x)\).  There are at most \(|T^*|=j\) such ratios.  Finally, if the internal radius is at least \(j/n\), then \((f,g)\) is column-close to \((c_f,c_g)\) and therefore cannot be CA-bad.
\end{proof}

\begin{proposition}[the boundary profile is Q]\label{prop:boundary-q}
Assume \(D\) is a multiplicative coset with domain polynomial \(X^n-\beta\).  Let
\[
        P_z(X)=X^m+\sum_{j=1}^w z_jX^{m-j},\qquad
        X^n=P_zQ_z+R_z,\quad \deg R_z<m,
\]
and put \(\omega=n-m\).  In the interpolation-lattice boundary branch of shifted degree \(w+1\), the prefix fiber is exactly the prescribed-divisor census
\[
        |\Fib_w(z)|
        =
        \#\{A:\deg A\le\omega-w-1,\ Q_z+A\mid X^n-\beta\}.
\]
Consequently Q is the boundary profile of the split-pencil problem, and BC starts only after this profile is removed.
\end{proposition}

\begin{proof}
The pair \((1,P_z)\) is in the interpolation module of \(U_z\).  The Euclidean division above gives the second vector \((Q_z,\beta-R_z)\), since for every \(x\in D\) one has \(Q_z(x)P_z(x)=\beta-R_z(x)\).  Their determinant is \(\beta-X^n\), so they form a basis.  In the boundary branch, monic elements of \(W\)-degree \(\omega\) are exactly \(Q_z+A\) with \(\deg A\le\omega-w-1\).  Being a valid locator is precisely the divisibility condition \(Q_z+A\mid X^n-\beta\).  This is the same set counted by \(\Fib_w(z)\), by \cref{prop:prefix-witness}.
\end{proof}

\begin{theorem}[interior base-field split-pencil floor]\label{prop:base-field-floor}
Let \(C=\RS[\F,D,K]\), where \(D\subseteq\B\subseteq\F\), \(|D|=n\), and put \(b=|\B|\).  Fix an agreement \(m\ge K\), put \(w=m-K\), and let
\[
        w+2\le d_1\le \left\lfloor \frac{n-K+1}{2}\right\rfloor .
\]
Set \(m_{d_1}=K-1+d_1\), and assume \(m\le m_{d_1}\le n\) and \(m_{d_1}+d_1\le n\).  Then there is a \(\B\)-valued received word \(U\) whose interpolation module has minimal shifted degree exactly \(d_1\), and whose size-\(m\) split-pencil census is at least
\[
        \binom{m_{d_1}}{m}
        \left\lceil\binom{n}{m_{d_1}}b^{-(d_1-1)}\right\rceil .
\]
Consequently any BC upper model for genuine interior profiles must include a base-field-normalized floor; a challenge-field-only normalization is false when \(|\F|\) is much larger than \(|\B|\).
\end{theorem}

\begin{proof}
For \(z=(z_1,\ldots,z_{d_1-1})\in\B^{d_1-1}\), let
\[
        U_z(X)=X^{m_{d_1}}+\sum_{h=1}^{d_1-1}z_hX^{m_{d_1}-h}
\]
on \(D\).  The module element \((1,U_z)\) has shifted degree \(m_{d_1}-(K-1)=d_1\).  If \((W,N)\in M_{U_z}\) had shifted degree \(<d_1\), then \(\deg W<d_1\), \(\deg N<m_{d_1}\), and \(WU_z-N\) would vanish on \(D\).  Its degree is at most \(m_{d_1}+d_1-1<n\), so \(WU_z=N\) as polynomials.  If \(W\ne0\), monicity of \(U_z\) gives \(\deg(WU_z)\ge m_{d_1}\), contradicting \(\deg N<m_{d_1}\); hence \(W=N=0\).  Thus the minimal shifted degree is exactly \(d_1\).

Now vary \(z\).  The first \(d_1-1\) high-degree coefficients of the locators of \(m_{d_1}\)-subsets of \(D\) lie in \(\B^{d_1-1}\), so by pigeonhole some \(z^*\) has at least
\[
        \left\lceil\binom n{m_{d_1}}b^{-(d_1-1)}\right\rceil
\]
subsets \(M'\) with that prefix.  For each such \(M'\), the polynomial \(U_{z^*}-\ell_{M'}\) has degree \(<K\), because the leading term and the first \(d_1-1\) coefficients cancel and \(m_{d_1}-d_1=K-1\).  Hence it agrees with \(U_{z^*}\) on every point of \(M'\).

Every \(m\)-subset \(T\subseteq M'\) is therefore an agreement support.  Distinct pairs \((M',T)\) give distinct codeword-support pairs: if the same degree-\(<K\) codeword agreed on the same \(T\) with \(|T|=m\ge K\), then the codewords coincide, and then \(\ell_{M'}\) coincides as well.  By \cref{prop:lattice-split}, these agreement-support pairs are represented in the split-pencil count of any shifted weak-Popov basis of \(M_{U_{z^*}}\).  Since the minimal shifted degree is \(d_1\), the family lies in the genuine interior profile \(d_1\), not in the boundary-Q profile.
\end{proof}

\begin{proposition}[rank-one pole-line base-field floor]\label{prop:rank-one-floor}
Let \(C=\RS[\F,D,K]\), with \(D\subseteq\B\subsetneq\F\), and fix \(m\ge K+1\), \(w=m-K\).  Let \(z^*\in\B^w\) be a heaviest depth-\(w\) prefix value among \(m\)-subsets of \(D\), and choose \(\alpha\in\F\setminus(D\cup\B)\).  For the pole line
\[
        f_\alpha(x)=\frac{U_{z^*}(x)}{x-\alpha},\qquad
        g_\alpha(x)=\frac{-1}{x-\alpha},
\]
there are at least
\[
        \left\lceil\binom{n}{m}|\B|^{-w}\right\rceil
\]
non-common rank-one agreement supports of size \(m\).  Moreover the line is genuinely extension-valued, even after common scalar multiplication.
\end{proposition}

\begin{proof}
By pigeonhole, \(|\Fib_w(z^*)|\ge\lceil\binom nm|\B|^{-w}\rceil\).  For each \(S\in\Fib_w(z^*)\), \cref{prop:pole-line} supplies the slope \(\zeta_S=U_{z^*}(\alpha)-\ell_S(\alpha)\) for which \(f_\alpha+\zeta_S g_\alpha\) is explained on \(S\).  In the slope-elimination notation of \cref{prop:slope-elimination}, this is a rank-one relation
\[
        (S_r(f_\alpha,S))_{r<w}+\zeta_S(S_r(g_\alpha,S))_{r<w}=0 .
\]
It is non-common because \(g_\alpha\) cannot be explained on any set of size at least \(K+1\), by the same argument as in \cref{prop:pole-line}: otherwise \((X-\alpha)G+1\) would have degree at most \(K\), vanish on at least \(K+1\) points, and take value \(1\) at \(\alpha\).

Finally, suppose some nonzero scalar \(\lambda\) made \(\lambda g_\alpha(D)\subseteq\B\).  For distinct \(x,y\in D\), the ratio \(g_\alpha(x)/g_\alpha(y)=(y-\alpha)/(x-\alpha)\) would lie in \(\B\).  Writing this ratio as \(r\ne1\) gives \(\alpha=(y-rx)/(1-r)\in\B\), a contradiction.
\end{proof}

\begin{proposition}[extension-valued distinct-slope rank-one floor]\label{prop:rank-one-distinct-slope-floor}
Let \(C=\RS[\F,D,K]\), with \(D\subseteq\B\subsetneq\F\), put \(q=|\F|\), and fix \(m\ge K+1\), \(w=m-K\).  Let \(z^*\in\B^w\) be a prefix value and let \(\mathcal G\subseteq\Fib_w(z^*)\) have size \(N\).  Then there exists \(\alpha\in\F\setminus\B\) such that the pole line
\[
        f_\alpha(x)=\frac{U_{z^*}(x)}{x-\alpha},
        \qquad
        g_\alpha(x)=\frac{-1}{x-\alpha}
\]
has at least
\[
        N-\frac{\binom N2 (K-1)}{q-|\B|}
\]
distinct finite MCA-bad slopes at agreement \(m\).  In particular, if \(N\le(q-|\B|)/(K-1)+1\), then some genuinely extension-valued pole line has at least \(N/2\) distinct MCA-bad slopes.  The same conclusion applies to any \(N\le\lceil\binom nm|\B|^{-w}\rceil\) by choosing \(\mathcal G\) inside a heaviest prefix fiber.
\end{proposition}

\begin{proof}
For each \(S\in\mathcal G\), \cref{prop:pole-line} gives the bad slope
\[
        \zeta_S(\alpha)=U_{z^*}(\alpha)-\ell_S(\alpha).
\]
If \(S\ne S'\), then \(\ell_S-\ell_{S'}\) has degree at most \(m-w-1=K-1\), because the leading term and the first \(w\) high-degree coefficients cancel.  Hence \(\zeta_S(\alpha)=\zeta_{S'}(\alpha)\) for at most \(K-1\) values of \(\alpha\in\F\), and therefore for at most \(K-1\) values in \(\F\setminus\B\).

Averaging over the \(q-|\B|\) allowed extension poles gives some \(\alpha\notin\B\) with at most \(\binom N2(K-1)/(q-|\B|)\) colliding support-pairs.  The number of distinct slope values is at least \(N\) minus that number of colliding pairs.  By \cref{prop:pole-line}, all these slopes are MCA-bad at agreement \(m\), because the pole pair is never simultaneously explained on any support of size at least \(K+1\).  Finally, since \(\alpha\notin\B\), the scalar-confinement argument from \cref{prop:rank-one-floor} shows that the line is genuinely extension-valued.
\end{proof}

\begin{remark}[proved reductions versus dense flatness]\label{rem:proved-not-flat}
The propositions and theorems in this section, together with the safe anchors above, are proved reductions, solved subcases, or route-cut counterexamples: prefix exactness, pole-line transport, fiber-to-slope conversion, exact high-agreement tangent cells, subfield confinement, the exponential aperiodic floor below the envelope, collision rigidity, second moments, twist-orbit moment amplification, quotient descent, coefficient-scale quotient extraction, the top-stratum quotient sieve, head-depth flatness, slope elimination, near-rational lattice classification, deficiency-one eliminants, identically split chart collapse to tangent cells, split pencils, the Q boundary profile, rank-one pole-line floors with distinct extension-valued slope realization, and base-field census floors.  They do not by themselves prove prefix flatness at the deployed frontier depth, the interior split-pencil census, or the finite adjacent safe inequalities.  They replace those questions by explicit fiber, slope, pencil, quotient, extension, and primitive shift-pair counts, and they prove the head-depth part of Q outright.
\end{remark}



\section{Proper Theorems for the Former Q/BC/SP Cells}

The labels Q, BC, and SP are retained only as names for residual safe-side mechanisms.  The section records actual theorems rather than certificate compilers.  The proved Q bound below is an unconditional packing ceiling, far weaker than the desired max-fiber theorem.  The BC material now has two parts: first, raw support locators are separated from the saturated codeword rays that MCA actually counts; second, the one-parameter primitive split-pencil obstruction is settled by a moving-root incidence bound.  Thus a residual BC chart that has genuinely been reduced to a projective locator pencil has only a constant number of split parameters in the deployed rows.  What remains is a chart-decomposition audit: every residual balanced-core chart must either be tangent/common/quotient/extension, or must present such a one-parameter moving-root pencil.  The proved SP statement has two forms: a crude all-pairs ceiling, and the sharper structural fact that SP is automatically controlled once Q is available as a max-fiber bound.  None of these statements is advertised as closing the dense frontier without the corresponding finite row ledger.

For the entropy calculations, let
\[
        \Xi_K(a)=\log_2\binom na-(a-K)\log_2|\B|,
\]
where \(K=k\) for a list route and \(K=k+1\) for the MCA lower route.  Also set
\[
        \varphi_\beta(g)=\Htwo(\rho+g)-\beta g,
        \qquad 0\le g\le1-\rho .
\]
The number \(g^*(\rho,\beta)\) is the last zero-crossing scale for \(\varphi_\beta\): \(\varphi_\beta(g)>0\) below the envelope and \(\varphi_\beta(g)<0\) above it.

\begin{lemma}[entropy bookkeeping for the frontier]\label{lem:entropy-bookkeeping}
Let \(K=k+O(1)\), \(k/n\to\rho\), and \(a=K+gn+O(1)\).  Then
\[
        \Xi_K(a)=n\varphi_\beta(g)+O(\log n+\beta).
\]
Consequently, if \(B_n^*\) is a target numerator and \(R(n)\ge1\) is a ledger overhead, then
\[
        R(n)2^{\Xi_K(a)}\le B_n^*
\]
holds eventually as soon as
\[
        n\bigl(-\varphi_\beta(g)\bigr)
        \ge
        \log_2R(n)-\log_2B_n^*+\omega(\log n+\beta).
\]
Similarly, the prefix floor exceeds the target, \(2^{\Xi_K(a)}>B_n^*\), as soon as
\[
        n\varphi_\beta(g)
        \ge
        \log_2B_n^*+\omega(\log n+\beta).
\]
If the residual factors and budgets are polynomial in \(n\), and if \(g^*=g^*(\rho,\beta)\) is the crossing point, then a logarithmic agreement reserve away from \(K+g^*n\) suffices.  More generally, for \(e^{o(n)}\) residual factors, any sublinear reserve whose entropy gain dominates their logarithms suffices.
\end{lemma}

\begin{proof}
Stirling's estimate gives, uniformly for \(a/n\) in compact subintervals of \((0,1)\),
\[
        \log_2\binom na=n\Htwo(a/n)+O(\log n).
\]
Since \(K=k+O(1)\) and \(a=K+gn+O(1)\), we have \(a/n=\rho+g+O(1/n)\) and
\[
        (a-K)\log_2|\B|=gn\beta+O(\beta).
\]
This proves the displayed asymptotic for \(\Xi_K(a)\).  The two integer inequalities follow by exponentiating after absorbing the \(O(\log n+\beta)\) term into the written \(\omega(\log n+\beta)\) margin.

Finally, for fixed \(\beta\), \(\varphi_\beta\) is strictly concave, \(\varphi_\beta(0)=\Htwo(\rho)>0\), and \(\varphi_\beta(1-\rho)=-\beta(1-\rho)<0\).  Hence the last crossing is isolated, and on either side of it the mean-value theorem converts an agreement reserve \(s\) into an entropy reserve comparable to \(|\varphi_\beta'(g^*)|s\), up to lower-order terms.  This is \(\Omega(s)\) for fixed \(\beta\), and \(\Omega(\beta s)\) in the large-base-field regime where the derivative is of order \(\beta\).  Thus \(s=O(\log n)\) handles polynomial overheads, while an \(o(n)\) reserve can be chosen to dominate any prescribed \(e^{o(n)}\) overhead.  If \(\beta=\beta_n\) varies with \(n\), the displayed entropy inequalities are the safe formulation; the phrase ``below the crossing'' must mean below the moving zero \(g^*(\rho,\beta_n)\) by enough entropy reserve to make those inequalities true.
\end{proof}

\begin{theorem}[identity-prefix unsafe theorem]\label{thm:unsafe-envelope}
Let \(D_n\subseteq\B_n\subseteq\F_n\), \(|D_n|=n\), \(k_n/n\to\rho\), \(\beta_n=\log_2|\B_n|\), and let \(B_n^*\) be the target numerator.  Let \(K_n=k_n+\eta\), where \(\eta=0\) for a list row and \(\eta=1\) for the MCA lower route, and let
\[
        a_n=K_n+g_nn+O(1),
        \qquad
        w_n=a_n-K_n .
\]
For list rows, if
\[
        n\varphi_{\beta_n}(g_n)
        \ge
        \log_2 B_n^*+\omega(\log n+\beta_n),
\]
then the list numerator at agreement \(a_n\) is larger than \(B_n^*\).  For MCA rows, put
\[
        L_n=\left\lceil \binom n{a_n}|\B_n|^{-w_n}\right\rceil,
        \qquad q_n=|\F_n| .
\]
If \(q_n>n\) and
\[
        \left\lceil \frac{L_n(q_n-n)}{q_n-n+k_n(L_n-1)}\right\rceil>B_n^*,
\]
then the MCA numerator at agreement \(a_n\) is larger than \(B_n^*\).  Therefore every agreement sequence below the relevant moving crossing, in the precise sense that the displayed entropy reserve holds for \(\varphi_{\beta_n}(g_n)\), is unsafe for list rows; it is unsafe for MCA rows whenever the exact simple-pole separation inequality above also holds.
\end{theorem}

\begin{proof}
By \cref{lem:entropy-bookkeeping}, the list hypothesis gives
\[
        \left\lceil\binom n{a_n}|\B_n|^{-w_n}\right\rceil>B_n^* .
\]
The list conclusion is exactly \cref{cor:identity-prefix-unsafe}(i).  For MCA, the same prefix list is built in \(\RS[\F_n,D_n,k_n+1]\); the displayed pole-separation inequality is exactly the hypothesis in \cref{cor:identity-prefix-unsafe}(ii).  That corollary gives more than \(B_n^*\) MCA-bad finite slopes at agreement \(a_n\).  Finally, if \(g_n\) stays below the last zero \(g^*(\rho,\beta_n)\) by enough reserve to make the displayed entropy inequality true, then the list conclusion follows.  This is deliberately stated in terms of \(\varphi_{\beta_n}\) so that fixed-base-field and growing-base-field rows use the same theorem without ambiguity.
\end{proof}

\begin{theorem}[Q: unconditional prefix-boundary packing bound]\label{thm:q-proper}
Assume \(1\le w<m\le n\).  Let \(\Phi_w:\binom Dm\to\B^w\) be the prefix map, and let \(\mathcal P_Q(z)\subseteq\Fib_w(z)\) be any pruned subfamily of a prefix fiber, for example the primitive part after quotient-pulled-back or planted members have been removed.  Put
\[
        t=\left\lfloor \frac{w}{2}\right\rfloor,
        \qquad
        V_t(n,m)=\sum_{i=0}^{t}\binom mi\binom{n-m}{i}.
\]
Then, for every prefix value \(z\),
\[
        |\mathcal P_Q(z)|\le |\Fib_w(z)|\le \frac{\binom nm}{V_t(n,m)}.
\]
Consequently the density-normalized prefix-boundary overhead satisfies
\[
        R_Q^{\rm pack}
        :=|\B|^w\max_z\frac{|\mathcal P_Q(z)|}{\binom nm}
        \le \frac{|\B|^w}{V_t(n,m)}
        \le \frac{|\B|^w}{\binom mt\binom{n-m}{t}}.
\]
The anticode cap
\[
        |\mathcal P_Q(z)|\le |\Fib_w(z)|\le \frac{\binom nm}{\binom{n-m+w}{w}}
\]
also holds.
\end{theorem}

\begin{proof}
Use the Johnson distance \(d_J(M,M')=|M\setminus M'|\) on \(m\)-subsets of \(D\).  By \cref{prop:prefix-rigidity}, two distinct members of a common prefix fiber have Johnson distance at least \(w+1\).  Therefore the closed Johnson balls of radius \(t=\lfloor w/2\rfloor\) around members of \(\Fib_w(z)\) are pairwise disjoint: if two such balls met, the two centers would have distance at most \(2t\le w\), a contradiction.

The radius-\(t\) Johnson ball around any \(m\)-subset has size \(V_t(n,m)\), because a point at distance \(i\) is obtained by deleting \(i\) elements of the center and inserting \(i\) elements from its complement.  The disjoint union of all these balls is contained in the full set \(\binom Dm\), which has size \(\binom nm\).  Hence \(|\Fib_w(z)|V_t(n,m)\le\binom nm\).  Since \(\mathcal P_Q(z)\subseteq\Fib_w(z)\), the same bound holds for \(\mathcal P_Q(z)\).  Multiplication by \(|\B|^w/\binom nm\) gives the normalized overhead inequality, and the final one-term bound follows from \(V_t(n,m)\ge\binom mt\binom{n-m}{t}\).  The anticode cap is \cref{cor:anticode-cap} rewritten using
\[
        \frac{\binom{n}{m-w}}{\binom mw}=\frac{\binom nm}{\binom{n-m+w}{w}}.
\]
\end{proof}

\begin{theorem}[BC: unconditional split-pencil dimension bound]\label{thm:bc-proper}
Fix \(C=\RS[\F,D,K]\).  Let \(U:D\to\F\), let \(m\ge K\), put \(\omega=n-m\), and let \(g_1,g_2\) be a shifted weak-Popov basis of the interpolation module \(M_U\) from \cref{prop:lattice-split} with shifted degrees \(d_1,d_2\).  Define
\[
        r_i=\max\{\omega-d_i+1,0\}\qquad(i=1,2).
\]
The number of size-\(m\) agreement supports represented in this split-pencil profile is at most
\[
        \min\left\{\binom n\omega, |\F|^{r_1+r_2}\right\}.
\]
If the basis and the allowed coefficient pencils are defined over \(\B\), and the cell requires \(A,B\in\B[X]\), the second term may be replaced by \(|\B|^{r_1+r_2}\).  For an MCA line, after common-support branches have been paid first, the number of finite bad slopes carried by this profile is bounded by the same quantity.
\end{theorem}

\begin{proof}
By \cref{prop:lattice-split}, every size-\(m\) agreement support in the profile is represented by
\[
        Ag_1+Bg_2,
        \qquad
        \deg A\le\omega-d_1,
        \qquad
        \deg B\le\omega-d_2,
\]
with the convention that no nonzero polynomial is allowed when the degree bound is negative.  The two coefficient spaces therefore have dimensions \(r_1\) and \(r_2\), so there are at most \(|\F|^{r_1+r_2}\) pairs \((A,B)\), and each pair determines at most one locator \(W=AW_1+BW_2\).  Independently, a monic \(D\)-split locator of degree \(\omega\) is the locator of an \(\omega\)-subset of \(D\), so there are at most \(\binom n\omega\) possible locators.  This proves the first displayed bound.  If the coefficient spaces are over \(\B\), the same dimension count gives \(|\B|^{r_1+r_2}\).

For the MCA assertion, \cref{prop:slope-elimination} says that a support that is not common to the two received words supports at most one finite slope.  The first-match convention assigns common-support branches before this split-pencil profile, so the remaining bad slopes are overcounted by the number of remaining locators.
\end{proof}

\begin{definition}[projective locator pencil and its common \(D\)-part]\label{def:projective-locator-pencil}
Let \(D\subseteq\F\) be finite, put \(\Lambda_D=\prod_{x\in D}(X-x)\), and let \(A,B\in\F[X]\) be not both zero.  The projective locator pencil generated by \(A,B\) is
\[
        L_{[s:t]}(X)=sA(X)+tB(X),\qquad [s:t]\in\mathbb P^1(\F).
\]
Its fixed \(D\)-part is
\[
        G_D(A,B)=\gcd(A,B,\Lambda_D),
        \qquad g_D(A,B)=\deg G_D(A,B).
\]
A root of \(G_D(A,B)\) is a common root of every member of the pencil; these are precisely the roots assigned to the common-GCD branch in the first-match ledger.
\end{definition}

\begin{theorem}[moving-root bound for one-parameter split pencils]\label{thm:bc-moving-root}
Let \(D\subseteq\F\) have size \(n\), and let \(L_{[s:t]}=sA+tB\) be a projective locator pencil as in \cref{def:projective-locator-pencil}.  Put \(G=G_D(A,B)\) and \(g=\deg G\).  Let \(\mathcal Z\subseteq\mathbb P^1(\F)\), and suppose that for every \(\lambda=[s:t]\in\mathcal Z\), the polynomial \(L_\lambda/G\) has at least \(h\ge1\) distinct roots in \(D\setminus Z(G)\).  Then
\[
        |\mathcal Z|\le \left\lfloor\frac{n-g}{h}\right\rfloor .
\]
In particular, if every counted member has degree \(j\), has fixed \(D\)-part exactly \(G\), and is \(D\)-split, then
\[
        |\mathcal Z|\le
        \left\lfloor\frac{n-g}{j-g}\right\rfloor,
        \qquad j>g.
\]
The case \(j=g\) is not a moving split-pencil cell: all \(D\)-roots are fixed and the member belongs to the common-GCD branch.
\end{theorem}

\begin{proof}
Count incidences
\[
        I=\{(\lambda,x):\lambda\in\mathcal Z,
        \ x\in D\setminus Z(G),\ L_\lambda(x)=0\}.
\]
For a fixed moving root \(x\in D\setminus Z(G)\), the pair \((A(x),B(x))\) is not \((0,0)\).  Hence the homogeneous linear equation
\[
        sA(x)+tB(x)=0
\]
has exactly one solution \([s:t]\in\mathbb P^1(\F)\).  Thus each moving domain point contributes to at most one incidence, and \(|I|\le n-g\).  On the other hand, every \(\lambda\in\mathcal Z\) contributes at least \(h\) incidences by hypothesis, so \(|I|\ge h|\mathcal Z|\).  This proves the first bound.

If a counted member is \(D\)-split of degree \(j\) and has fixed \(D\)-part exactly \(G\), then exactly \(j-g\) of its \(D\)-roots are moving roots in \(D\setminus Z(G)\).  Applying the first bound with \(h=j-g\) gives the displayed inequality.  If \(j=g\), no moving root remains; the split condition has already been paid by the common-GCD branch rather than by a primitive moving pencil.
\end{proof}

\begin{corollary}[BC is settled on primitive one-parameter locator pencils]\label{cor:bc-one-pencil}
Consider an MCA split-pencil chart at agreement \(m\), and write \(\omega=n-m\).  Suppose that after the first-match tangent, common-support, quotient, extension, degree-drop, and common-GCD branches have been removed, the residual candidate error locators in this chart are represented by a projective locator pencil \(L_\lambda=sA+tB\), with the slope parameter injecting into \(\lambda\), and that every residual bad slope in the chart has at least one degree-\(\omega\) \(D\)-split locator in the pencil.  Then the number of residual bad finite slopes in this chart is at most
\[
        \left\lfloor\frac{n-g}{\omega-g}\right\rfloor,
        \qquad g=g_D(A,B)<\omega .
\]
In the primitive case \(g=0\), this is \(\lfloor n/(n-m)\rfloor\).  For the two active MCA adjacent rows in \cref{rem:finite-artifacts},
\[
\begin{array}{c|c|c|c}
\text{row} & m=a_+ & \omega=n-m & \lfloor n/\omega\rfloor\\ \hline
\text{KoalaBear MCA} & 1116048 & 980104 & 2\\
\text{Mersenne-31 MCA} & 1116024 & 980128 & 2
\end{array}
\]
so every primitive one-parameter BC pencil contributes at most two finite bad slopes.
\end{corollary}

\begin{proof}
The first-match removals ensure that the chart is counted only through moving roots of the displayed pencil.  A residual bad slope supplies at least one degree-\(\omega\) split locator, and the slope-to-pencil parameter is injective by hypothesis; hence the slope count is bounded by the number of split parameters.  Apply \cref{thm:bc-moving-root} with \(j=\omega\).  The displayed row values use \(n=2^{21}\), \(a_+=1116048\) for KoalaBear MCA, and \(a_+=1116024\) for Mersenne-31 MCA.
\end{proof}

\begin{remark}[what this does and does not settle for BC]\label{rem:bc-status-after-moving-root}
\Cref{thm:bc-moving-root,cor:bc-one-pencil} settle the enumerative part of BC for the object named in the original program: a primitive one-parameter split-locator pencil.  The proof counts slopes, not raw supports, and it automatically incorporates the saturation correction because multiple size-\(m\) supports below one saturated codeword ray do not create new pencil parameters.  Therefore a BC chart that has been reduced to such a pencil is no longer a conjectural input.

The remaining BC audit is structural rather than enumerative.  For a finite deployed row one must verify that each residual balanced-core chart produced by the interpolation-lattice reduction is one of the paid types, or else supplies a projective locator pencil to which \cref{cor:bc-one-pencil} applies.  A higher-dimensional coefficient family not reducible to one-parameter moving-root pencils is not covered by this theorem and must be split further or charged to a separate named cell.
\end{remark}


\begin{definition}[saturated codeword rays and support census]\label{def:saturated-rays}
Let \(C=\RS[\F,D,K]\), let \(U:D\to\F\), and put
\[
        S_c(U)=\{x\in D:U(x)=c(x)\},\qquad s_c(U)=|S_c(U)|
\]
for \(c\in C\).  The saturated ray set at agreement \(m\) is
\[
        \operatorname{Ray}(U;m)=\{c\in C:s_c(U)\ge m\}.
\]
Let \(\omega=n-m\), and define the support-locator census
\[
\begin{aligned}
\Cen(U;m)=\#\{(W,N)\in \F[X]^2:\;& W\text{ monic},\ \deg W=\omega,\ W\mid \ell_D,\\
        & WU=N\text{ on }D,\quad W\mid N,\quad \deg(N/W)<K\}.
\end{aligned}
\]
Equivalently, \(W=\ell_{D\setminus T}\) for an agreement support \(T\subseteq D\) of size \(m\).
\end{definition}

\begin{theorem}[exact saturation identity]\label{thm:saturation}
With \(C=\RS[\F,D,K]\) and \(\Cen(U;m)\) as in \cref{def:saturated-rays}, for every word \(U:D\to\F\),
\[
        \Cen(U;m)=\sum_{c\in C}\binom{s_c(U)}m .
\]
Moreover, the support locators lying above a fixed codeword ray \(c\) are exactly
\[
        (W,N)=H\,(G_c,G_cc),
        \qquad
        G_c=\ell_{D\setminus S_c(U)},
        \qquad
        H\mid\ell_{S_c(U)},
        \qquad
        \deg H=s_c(U)-m .
\]
Consequently the fiber size of the forgetful map from support locators to saturated codeword rays is \(\binom{s_c(U)}m\) above \(c\).
\end{theorem}

\begin{proof}
Let \((W,N)\) be counted by \(\Cen(U;m)\).  Since \(W\) is monic, divides \(\ell_D\), and has degree \(n-m\), there is a unique \(m\)-set \(T\subseteq D\) such that \(W=\ell_{D\setminus T}\).  Since \(W\mid N\) and \(\deg(N/W)<K\), write \(N=Wc\) for a unique codeword \(c\in C\).  On every \(x\in T\), the value \(W(x)\) is nonzero; the identity \(WU=N=Wc\) on \(D\) therefore gives \(U(x)=c(x)\).  Thus \(T\subseteq S_c(U)\).

Conversely, if \(c\in C\) and \(T\subseteq S_c(U)\) has \(|T|=m\), then \(W=\ell_{D\setminus T}\) and \(N=Wc\) satisfy all defining conditions of \(\Cen(U;m)\).  For fixed \(c\), the number of such supports is \(\binom{s_c(U)}m\), and summing over \(c\) gives the first identity.

For the fiber statement, write \(G_c=\ell_{D\setminus S_c(U)}\).  If \(T\subseteq S_c(U)\), then
\[
        \ell_{D\setminus T}
        =\ell_{D\setminus S_c(U)}\ell_{S_c(U)\setminus T}
        =G_cH,
\]
where \(H=\ell_{S_c(U)\setminus T}\), hence \(H\mid\ell_{S_c(U)}\) and \(\deg H=s_c(U)-m\).  The converse is the same construction reversed, because every monic divisor \(H\) of \(\ell_{S_c(U)}\) of degree \(s_c(U)-m\) is the locator of \(S_c(U)\setminus T\) for a unique \(m\)-set \(T\subseteq S_c(U)\).
\end{proof}

\begin{corollary}[raw support BC is not the MCA object]\label{cor:raw-bc-fails}
In the setting of \cref{thm:saturation}, if \(U\) agrees with a single codeword \(c\) on \(n-d\) positions and with no other codeword on \(m\) or more positions, then
\[
        |\operatorname{Ray}(U;m)|=1,
        \qquad
        \Cen(U;m)=\binom{n-d}{m} .
\]
Thus an exponentially large raw support-locator family may represent one MCA ray.  A split-pencil proof for MCA must therefore count saturated primitive line-rays after quotient, prefix, tangent, and common-support branches have been paid; a raw support census is a list/moment object, not the final MCA numerator.
\end{corollary}

\begin{proof}
The ray statement is the hypothesis.  The support-census identity is \cref{thm:saturation} with one nonzero summand \(\binom{s_c(U)}m=\binom{n-d}m\).  The final sentence is just the same equality interpreted in MCA terms: MCA counts a bad slope once, whereas \(\Cen(U;m)\) counts all size-\(m\) supports below the same saturated agreement ray.
\end{proof}

\begin{definition}[line rays]\label{def:line-rays}
In the setting of \cref{def:saturated-rays}, for an affine line \(U_Z=u+Zv\) and a set of finite slopes \(E\subseteq\F\), define
\[
        \LineRay_E(u,v;m)
        =
        \{(z,c)\in E\times C:\ |\{x\in D:u(x)+zv(x)=c(x)\}|\ge m\}.
\]
Write \(s_{z,c}=|\{x\in D:u(x)+zv(x)=c(x)\}|\).
\end{definition}

\begin{proposition}[line-ray saturation identity]\label{prop:line-ray-saturation}
With notation as in \cref{def:line-rays}, for every \(E\subseteq\F\),
\[
        \sum_{z\in E}\Cen(U_z;m)
        =
        \sum_{(z,c)\in\LineRay_E(u,v;m)}\binom{s_{z,c}}m .
\]
Consequently
\[
        |\{z\in E:\exists c\ (z,c)\in\LineRay_E(u,v;m)\}|
        \le
        |\LineRay_E(u,v;m)|
        \le
        \sum_{z\in E}\Cen(U_z;m).
\]
The two inequalities are the exact deduplication losses between raw support locators, saturated codeword rays, and slope numerators.
\end{proposition}

\begin{proof}
Apply \cref{thm:saturation} separately to the word \(U_z=u+zv\) for each \(z\in E\), and then sum over \(z\).  The first inequality forgets the codeword coordinate of a line ray, and the second follows because each line ray has \(s_{z,c}\ge m\), hence contributes at least one support to the right-hand side of the identity.  \Cref{cor:raw-bc-fails} shows that the second inequality can be exponentially loose for a single slope and a single codeword ray, so an MCA upper ledger must either count slopes directly or control this saturation factor.
\end{proof}

\begin{theorem}[Q discharges the shift-pair ledger]\label{thm:q-implies-sp}
Let \(N_w(z)=|\Fib_w(z)|\), let
\[
        \overline F=\binom nm |\B|^{-w},
\]
and let \(P_e\) denote the full off-diagonal shift-pair stratum in \cref{prop:second-moment}, namely
\[
        P_e=\sum_{\substack{R\subseteq D\\ |R|=m-e}}\operatorname{sp}_w(e;D\setminus R).
\]
If Q holds in max-fiber form with constant \(\kappa\),
\[
        \max_z N_w(z)\le \kappa\overline F,
\]
then necessarily \(\kappa\overline F\ge1\) whenever \(\binom Dm\) is nonempty, and
\[
        \sum_{e=w+1}^{\min(m,n-m)}P_e
        \le
        \left(\kappa-\frac{1}{\overline F}\right)\binom nm\,\overline F .
\]
Every quotient-removed or primitive shift-pair subledger satisfies the same upper bound.  Hence SP is not an independent residual input in any closing theorem that already proves Q as a max-fiber theorem.
\end{theorem}

\begin{proof}
The exact second-moment identity of \cref{prop:second-moment} says
\[
        \sum_z N_w(z)^2
        =\binom nm+
        \sum_{e=w+1}^{\min(m,n-m)}P_e .
\]
On the other hand,
\[
        \sum_zN_w(z)^2
        \le (\max_zN_w(z))\sum_zN_w(z)
        \le \kappa\overline F\binom nm,
\]
because \(\sum_zN_w(z)=\binom nm\).  Since at least one fiber is nonempty when \(\binom Dm\ne\varnothing\), the hypothesis also implies \(\kappa\overline F\ge1\), so the displayed right-hand side is nonnegative in the nonvacuous case.  Subtracting the diagonal term \(\binom nm\) gives the displayed inequality.  Any primitive or quotient-removed ledger is a sub-sum of the full off-diagonal ledger, so it is bounded a fortiori.
\end{proof}

\begin{theorem}[SP: unconditional primitive shift-pair bound]\label{thm:sp-proper}
Use the notation of \cref{prop:gamma2-ledger}.  For \(w+1\le e\le\min(m,n-m)\), put
\[
        T_e(n,m)=\binom n{m-e}\binom{n-m+e}{e}\binom{n-m}{e}.
\]
Then the total off-diagonal contribution at distance \(e\) satisfies
\[
        P_e^{\rm quot}+P_e^{\rm prim}\le T_e(n,m).
\]
If \(D\subseteq\B\), then the top stratum \(e=w+1\) also satisfies the sharper constant-shift bound
\[
        P_{w+1}^{\rm quot}+P_{w+1}^{\rm prim}
        \le
        (|\B|-1)\binom n{m-w-1}\binom{n-m+w+1}{w+1}.
\]
Consequently
\[
        \Gamma_2
        \le
        \frac{1}{\overline F}
        +\frac{\sum_{e=w+1}^{\min(m,n-m)}T_e(n,m)}{\binom{n}{m}\overline F},
\]
with the option of replacing the \(e=w+1\) term by the sharper displayed top-stratum bound.
\end{theorem}

\begin{proof}
Fix \(e\).  In the exact second-moment decomposition of \cref{prop:second-moment}, an off-diagonal pair is encoded by a common part \(R\subseteq D\) of size \(m-e\), then by two disjoint degree-\(e\) root sets inside \(D\setminus R\).  There are \(\binom n{m-e}\) choices for \(R\).  Once \(R\) is fixed, the ambient set \(D\setminus R\) has size \(n-m+e\); the first root set has at most \(\binom{n-m+e}{e}\) choices and the second, disjoint from it, has at most \(\binom{n-m}{e}\) choices.  This ignores the depth equation and the quotient/primitive first-match assignment, so it is an upper bound for \(P_e^{\rm quot}+P_e^{\rm prim}\).

When \(e=w+1\), \cref{prop:second-moment} shows that every top-stratum shift pair has the form \(A-B=c\in\F^\times\).  If both polynomials split over \(D\subseteq\B\), then both have coefficients in \(\B\), hence \(c\in\B^\times\).  For each fixed \(R\), choose \(A\) by choosing its \(e\)-element root set in \(D\setminus R\), and choose \(c\in\B^\times\); the polynomial \(B=A-c\) is then determined.  Requiring \(B\) to split over \(D\setminus R\), to be disjoint from \(A\), and then assigning the pair to either the quotient or primitive ledger can only reduce the count.  This gives the top-stratum bound.  The displayed \(\Gamma_2\) inequality is \cref{prop:gamma2-ledger} with the total quotient-plus-primitive off-diagonal contribution in every stratum replaced by the just-proved upper bound.
\end{proof}

\begin{proposition}[size of the proper Q bound at the active adjacent rows]\label{prop:proper-q-gap}
At the four active adjacent agreements, the packing estimate in \cref{thm:q-proper} gives the following one-term upper bounds for the remaining normalized prefix overhead:
\[
\begin{array}{c|ccc}
\text{row} & w & \log_2 R_Q^{\rm pack}\text{ allowed by the bound} & \text{spare margin}\\
\hline
\text{KoalaBear MCA} & 67471 & 1661552.48 & 22.1969\\
\text{KoalaBear list} & 67471 & 1661552.47 & 22.0109\\
\text{Mersenne-31 MCA} & 67447 & 1660926.12 & 3.2589\\
\text{Mersenne-31 list} & 67447 & 1660926.11 & 3.0730 .
\end{array}
\]
Thus the unconditional packing proof is a real theorem but is far too weak to supply the adjacent safe ledger.
\end{proposition}

\begin{proof}
For each row, \(n=2097152\), \(K=k+1=1048577\) on the MCA route, \(K=k=1048576\) on the list route, \(w=a_+-K\), and \(\log_2|\B|=31\).  The last inequality of \cref{thm:q-proper} gives
\[
        \log_2 R_Q^{\rm pack}
        \le
        31w-\log_2\binom{a_+}{\lfloor w/2\rfloor}
             -\log_2\binom{n-a_+}{\lfloor w/2\rfloor}.
\]
Substituting the four active adjacent agreements \(1116048,1116047,1116024,1116023\) gives the displayed values.  For these four rows the Johnson-ball summand
\[
        J_i=\binom{a_+}{i}\binom{n-a_+}{i}
\]
is increasing on \(0\le i\le t\).  Indeed
\[
        \frac{J_i}{J_{i-1}}
        =
        \frac{(a_+-i+1)(n-a_+-i+1)}{i^2},
\]
and at \(i=t\) this ratio is already \(>900\) in all four rows, while the ratio is larger for smaller \(i\).  Hence \(V_t(n,a_+)\le(t+1)J_t\), so replacing the one-term denominator by the full ball can lower the displayed logarithmic upper bound by less than \(\log_2(t+1)<16\) bits.  This is negligible compared with the gap between the displayed million-bit overhead and the row margins.
\end{proof}

\begin{theorem}[theorem-level inventory with proper bounds]\label{thm:proper-inventory}
Without any dense-frontier census assumption, the following theorem-level pieces are available in this note.
\begin{enumerate}[label=\textup{(\roman*)}]
\item Prefix witnesses and the simple-pole list-to-MCA conversion give the exact lower-side criteria of \cref{prop:prefix-witness,thm:simple-pole-list-floor,cor:identity-prefix-unsafe,thm:unsafe-envelope}.
\item The broad safe side contains the deep-regime MCA theorem, the exact high-agreement tangent cell, MCA-from-CA below half distance, and subfield confinement, namely \cref{thm:gf-deep-mca,prop:exact-tangent-cell,thm:gf-mca-from-ca,thm:subfield-confinement}.
\item The residual cells have unconditional theorem-level pieces: Q has the Johnson packing and anticode bounds of \cref{thm:q-proper}; raw split-pencil support counts have the dimension bound of \cref{thm:bc-proper}; primitive one-parameter BC pencils are bounded by \cref{thm:bc-moving-root,cor:bc-one-pencil}; MCA support multiplicities are exactly saturated by \cref{thm:saturation,cor:raw-bc-fails,prop:line-ray-saturation}; and SP is downstream from Q by \cref{thm:q-implies-sp}, with the all-pairs ceiling of \cref{thm:sp-proper} available as an audit bound.
\item The stronger dense-frontier conclusion is not proved by these bounds.  In particular, \cref{prop:proper-q-gap} shows that the proved Q packing bound alone leaves overhead far beyond the active adjacent margins.  On the BC side, \cref{cor:raw-bc-fails} explains why raw support counts are the wrong numerator, while \cref{thm:bc-moving-root,cor:bc-one-pencil} settles the primitive one-parameter pencil count once the finite chart decomposition has supplied such a pencil.
\end{enumerate}
\end{theorem}

\begin{proof}
Parts (i) and (ii) are the cited theorems.  Part (iii) is exactly \cref{thm:q-proper,thm:bc-proper,thm:bc-moving-root,cor:bc-one-pencil,thm:saturation,cor:raw-bc-fails,prop:line-ray-saturation,thm:q-implies-sp,thm:sp-proper}.  Part (iv) follows from \cref{prop:proper-q-gap}: an upper estimate that permits more than \(2^{1.66\cdot10^6}\) normalized prefix overhead cannot imply a ledger with only \(2^{22.2}\) or less spare overhead.  The saturation identity changes the MCA split-pencil target from raw supports to primitive line-rays, and the moving-root theorem then pays the one-parameter primitive pencil case.  Therefore the adjacent safe inequality now requires stronger Q and a finite verification that the residual BC charts have all reduced to the paid moving-root form.
\end{proof}

\begin{theorem}[deployed-row theorem-level status]\label{thm:deployed-row-status}
For the four deployed rows and budgets displayed in \cref{rem:finite-artifacts}, the following unconditional conclusions follow from the banked integer arithmetic and the proper bounds in this note.
\begin{enumerate}[label=\textup{(\roman*)}]
\item The active agreements \(a_0\) in the row table are unsafe.
\item The identity-prefix lower route stops at \(a_0+1\); this is not an upper bound.
\item The broad safe anchors listed in \cref{prop:finite-safe-anchor-gap} start strictly above \(a_0+1\), so they do not prove the adjacent safe inequality.
\item The unconditional Q packing bound, raw split-pencil dimension bound, moving-root bound for one-parameter BC pencils, saturation identity, and SP-from-Q theorem of \cref{thm:q-proper,thm:bc-proper,thm:bc-moving-root,cor:bc-one-pencil,thm:saturation,thm:q-implies-sp} are proved.  They do not by themselves imply \(B(a_0+1)\le B^*\), because Q is still too weak and the finite BC chart-decomposition audit has not been supplied in this note.
\end{enumerate}
\end{theorem}

\begin{proof}
Part (i) and the lower-route stopping statement in part (ii) are \cref{prop:finite-packet-consequences}.  Part (iii) is \cref{prop:finite-safe-anchor-gap}.  Part (iv) follows because the first residual replacement bound already has the gap quantified in \cref{prop:proper-q-gap}; \cref{thm:bc-moving-root,cor:bc-one-pencil} pays only those BC charts that the finite reduction has already turned into one-parameter primitive locator pencils; and \cref{thm:q-implies-sp} only discharges SP after a sufficiently sharp Q max-fiber theorem has been proved.  Hence the note still does not supply the missing inequality \(B(a_0+1)\le B^*\).
\end{proof}

\begin{theorem}[adjacent criterion from an explicit upper bound]\label{thm:finite}
Fix one of the four deployed rows in the table, with its relevant numerator \(B(a)\) from \cref{def:staircase}.  Suppose the unsafe inequality at \(a_0\) is the banked identity-prefix theorem, and suppose a proved upper bound gives
\[
        B(a_0+1)\le U(a_0+1)\le B^* .
\]
For an MCA row, \(U\) may be obtained by a first-match, deduplicated sum of tangent, quotient, extension, prefix-boundary, split-pencil-interior, primitive shift-pair, and any explicitly named residual cells.  For a list row, \(U\) is the corresponding exact list upper bound.  Then the first safe integer agreement is \(a_0+1\), and \(\delta_C^*(\varepsilon^*)\) is pinned to one integer step under the closed-ball convention for that row object.
\end{theorem}

\begin{proof}
By the unsafe theorem, \(B(a_0)>B^*\), so \cref{lem:integer-budget} gives error \(>\varepsilon^*\) at agreement \(a_0\).  The proved upper bound gives \(B(a_0+1)\le U(a_0+1)\le B^*\).  Applying \cref{lem:integer-budget} again gives error at most \(\varepsilon^*\) at agreement \(a_0+1\).  Monotonicity makes every agreement threshold \(a\le a_0\) unsafe once \(a_0\) is unsafe, so \(a_0+1\) is the first safe integer agreement.  This is precisely the adjacent staircase condition of \cref{def:staircase}.  In the MCA case, \cref{lem:first-match-ledger} is the generic mechanism by which a first-match cell ledger proves the required upper inequality.
\end{proof}

\begin{remark}[endpoint convention]\label{rem:endpoint}
If \cref{thm:finite} is instantiated with a proved upper bound at \(a_0+1\), then the largest safe closed-grid radius is \((n-a_0-1)/n\).  The first proved unsafe grid radius is \((n-a_0)/n\), and monotonicity makes every larger closed-grid radius unsafe.  For MCA this is \cref{lem:basic-staircase}; for list rows it is the immediate fact that a codeword list at agreement \(a\) is still valid at every lower agreement threshold.  If one records \(\delta_C^*(\varepsilon^*)\) as a real supremum over closed integer balls, the threshold lies at the upper edge of this unit interval and the convention must be printed with the proof.
\end{remark}


\section{Attempted Q Closure: Exact Atom Target and Barrier}

This section records the result of trying to close the remaining prefix-boundary cell.  The useful outcome is an exact finite target and several no-shortcut checks.  The outcome is not a proof of the adjacent safe rows: the atom bound in \cref{def:q-row-atom} is precisely the missing theorem.

\begin{definition}[row-sharp Q atom bound]\label{def:q-row-atom}
Fix a deployed adjacent row, let \(a_+\) be its adjacent agreement, let \(K\) be the list dimension used on that route, and put \(w=a_+-K\).  For a first-match residual prefix-boundary family \(\mathcal P_Q(z)\subseteq\Fib_w(z)\), define
\[
        R_Q^{\max}
        =
        |\B|^w\max_z\frac{|\mathcal P_Q(z)|}{\binom n{a_+}}.
\]
The row-sharp Q atom bound with bit margin \(\Delta_Q\) is
\[
        R_Q^{\max}\le 2^{\Delta_Q}.
\]
Equivalently, if the entire row budget is hypothetically assigned to Q, the sufficient integer form is
\[
        \max_z|\mathcal P_Q(z)|\le B^* .
\]
Any non-Q cell paid in the same first-match ledger only decreases the available \(\Delta_Q\).
\end{definition}

\begin{proposition}[exact finite Q atom target at the four adjacent rows]\label{prop:q-exact-target}
At the active adjacent agreements, the full-budget Q atom target is as follows:
\[
\begin{array}{c|r|r|r|r}
\text{row} & w & \left\lceil \binom n{a_+}|\B|^{-w}\right\rceil & B^* & B^*/\left\lceil \binom n{a_+}|\B|^{-w}\right\rceil\\
\hline
\text{KoalaBear MCA} & 67471 & 57198030366 & 274980728111395087 & 4807520.9295\\
\text{KoalaBear list} & 67471 & 65065153468 & 274980728111395087 & 4226236.5253\\
\text{Mersenne-31 MCA} & 67447 & 1752700 & 16777215 & 9.5722\\
\text{Mersenne-31 list} & 67447 & 1993678 & 16777215 & 8.4152
\end{array}
\]
In bits, these are exactly the spare margins already printed in the row table:
\[
        22.1969,
        \quad 22.0109,
        \quad 3.2589,
        \quad 3.0730 .
\]
Thus the binding Q problem is the Mersenne-31 list row: after all first-match payments, the remaining primitive prefix-boundary atom must be at most about \(8.42\) times its average size, and less if any other cell consumes budget.
\end{proposition}

\begin{proof}
For each row, substitute \(n=2^{21}\), \(k=2^{20}\), the displayed adjacent agreement \(a_+\), and \(K=k+1\) on the MCA route and \(K=k\) on the list route.  The average prefix fiber is \(\binom n{a_+}|\B|^{-w}\), and the integer lower-route replay uses its ceiling.  The budgets are
\[
        B^*_{\rm KB}=\left\lfloor\frac{(2^{31}-2^{24}+1)^6}{2^{128}}\right\rfloor,
        \qquad
        B^*_{\rm M31}=\left\lfloor\frac{(2^{31}-1)^4}{2^{100}}\right\rfloor .
\]
The displayed ratios are the exact integer quotients evaluated as decimals.  Taking base-two logarithms of \(B^*/(\binom n{a_+}|\B|^{-w})\) gives the four printed bit margins.
\end{proof}

\begin{proposition}[moment-order floor for a finite Q proof]\label{prop:q-moment-order-floor}
Any proof that uses only the moment criterion of \cref{thm:moment-q} and a nontrivial bound on \(\Gamma_r\) must use moment order at least
\[
        r\ge\left\lceil\frac{w\log_2|\B|}{\Delta_Q}\right\rceil .
\]
For the active rows this lower floor is
\[
\begin{array}{c|r|r|r}
\text{row} & w & \Delta_Q\text{ bits} & \text{minimum }r\\
\hline
\text{KoalaBear MCA} & 67471 & 22.1969 & 94196\\
\text{KoalaBear list} & 67471 & 22.0109 & 94991\\
\text{Mersenne-31 MCA} & 67447 & 3.2589 & 641593\\
\text{Mersenne-31 list} & 67447 & 3.0730 & 680397
\end{array}
\]
Consequently every fixed-moment Q argument, and every asymptotic argument that does not produce moment estimates at orders comparable with \(w/\Delta_Q\), is numerically incapable of proving the finite adjacent rows.
\end{proposition}

\begin{proof}
For the normalized moments in \cref{prop:moment-sandwich}, Holder's inequality gives \(\Gamma_r\ge1\).  Indeed, the sum of the \(\mu(z)\)'s is one and there are at most \(|\B|^w\) possible prefix values, so the uniform distribution minimizes \(\sum_z\mu(z)^r\), giving \(\Gamma_r\ge1\).  The finite criterion in \cref{thm:moment-q} requires
\[
        \log_2\Gamma_r+w\log_2|\B|\le r\Delta_Q .
\]
Dropping the nonnegative term \(\log_2\Gamma_r\) gives the necessary inequality \(r\ge w\log_2|\B|/\Delta_Q\).  Here the table uses the exact row prime in \(\log_2|\B|\); for Mersenne-31 this differs from \(31\) only below the displayed precision, while for KoalaBear it is slightly below \(31\).
\end{proof}

\begin{proposition}[the BC moving-root theorem does not pay boundary Q]\label{prop:bc-not-q}
In the boundary-Q profile of \cref{prop:boundary-q}, the unknown divisor is
\[
        Q_z+A,
        \qquad
        \deg A\le \omega-w-1,
        \qquad
        \omega=n-a_+ .
\]
Thus the coefficient parameter space has dimension \(\omega-w\).  At the active rows this dimension is
\[
\begin{array}{c|r|r|r}
\text{row} & \omega=n-a_+ & w & \omega-w\\
\hline
\text{KoalaBear MCA} & 981104 & 67471 & 913633\\
\text{KoalaBear list} & 981105 & 67471 & 913634\\
\text{Mersenne-31 MCA} & 981128 & 67447 & 913681\\
\text{Mersenne-31 list} & 981129 & 67447 & 913682
\end{array}
\]
Therefore \cref{thm:bc-moving-root,cor:bc-one-pencil} cannot be applied directly to boundary Q: those theorems count split parameters on a projective one-parameter locator pencil, while Q is a high-dimensional affine divisor slice.  To route Q through BC one would need an additional theorem decomposing this affine slice into paid one-parameter pencils with a first-match bound on the number of relevant pencils.  No such theorem is present in this note.
\end{proposition}

\begin{proof}
The boundary dictionary \cref{prop:boundary-q} states exactly that a prefix fiber is counted by the divisibility condition \(Q_z+A\mid X^n-\beta\) with \(\deg A\le\omega-w-1\).  A polynomial of degree at most \(\omega-w-1\) has \(\omega-w\) free coefficients.  Substituting the four values of \(a_+\), \(K\), and \(w=a_+-K\) gives the table.  Since \(\omega-w\) is larger than nine hundred thousand in all four cases, the family is not a one-parameter pencil.  The moving-root theorem is still valid for every line inside this affine space, but a line-by-line decomposition without a bound on the number of lines gives no row budget.
\end{proof}

\begin{proposition}[current Q ingredients do not close the adjacent ledger]\label{prop:q-not-closed}
The theorem-level Q material currently proved in this note does not imply \(B(a_0+1)\le B^*\) for any deployed row.
\end{proposition}

\begin{proof}
The only unconditional max-fiber upper bound presently proved for Q is the Johnson packing/anticode bound of \cref{thm:q-proper}.  By \cref{prop:proper-q-gap}, its normalized overhead at the active rows is between \(2^{1660926}\) and \(2^{1661553}\), whereas the available row margins are only \(2^{22.1969}\), \(2^{22.0109}\), \(2^{3.2589}\), and \(2^{3.0730}\).  Hence that theorem is many orders of magnitude too weak.

\Cref{thm:moment-q} could in principle prove the desired atom bound, but \cref{prop:q-moment-order-floor} shows that the finite rows require moment orders from \(94196\) up to \(680397\) before the harmless term \(w\log_2|\B|\) alone fits in the margin.  No theorem in the note supplies such high-order moment estimates.

Finally, \cref{thm:q-implies-sp} is one-way: it proves that a max-fiber Q theorem would discharge SP, not that SP proves Q.  And \cref{prop:bc-not-q} explains why the one-parameter BC moving-root theorem does not pay the high-dimensional boundary-Q slice.  Thus the adjacent upper ledger remains open precisely at the row-sharp Q atom bound of \cref{def:q-row-atom}.
\end{proof}

\section{Audit Corrections and Expanded Gap Ledger}

The preceding sections now leave three kinds of statements, and the distinction is part of the proof.

\begin{proposition}[closed audit of what is and is not proved]\label{prop:closed-audit}
In this expanded version:
\begin{enumerate}[label=\textup{(\roman*)}]
\item The lower-side identity-prefix and simple-pole implications are proved theorems under their stated field-size hypotheses.
\item The broad safe anchors and the high-agreement tangent/common-line cells are proved only in their displayed agreement ranges.
\item The Q, BC, and SP replacements in \cref{thm:q-proper,thm:bc-proper,thm:bc-moving-root,cor:bc-one-pencil,thm:saturation,prop:line-ray-saturation,thm:q-implies-sp,thm:sp-proper} are unconditional statements.  The one-parameter primitive BC pencil is now numerically strong, but the full adjacent row is still not closed because the row ledger must verify that all residual BC charts reduce to such pencils and because Q remains too weak.
\item No statement in this note proves \(B(a_0+1)\le B^*\) for any deployed adjacent row.
\end{enumerate}
\end{proposition}

\begin{proof}
Part (i) is \cref{prop:prefix-witness,thm:simple-pole-list-floor,cor:identity-prefix-unsafe,thm:unsafe-envelope}.  Part (ii) is \cref{thm:gf-deep-mca,prop:exact-tangent-cell,thm:gf-mca-from-ca,thm:subfield-confinement,prop:finite-safe-anchor-gap}.  Part (iii) follows from the cited residual-cell theorems, from the moving-root count for one-parameter BC pencils, and from \cref{prop:proper-q-gap}.  Part (iv) is then immediate from \cref{thm:deployed-row-status}: the only theorem in this note that would turn the banked unsafe edge into a closed adjacent threshold is \cref{thm:finite}, and its upper-bound hypothesis is not supplied by the proved replacement bounds.
\end{proof}

\begin{problem}[row-sharp multilevel Q target]\label{prob:row-sharp-q}
For each active adjacent row, after quotient-pulled-back, planted, tangent, and explicitly paid lower-dimensional cells have been assigned by first match, prove a max-fiber bound of the form
\[
        \max_z |\mathcal P_Q(z)|
        \le
        R_Q^{\rm row}\binom n{a_+}|\B|^{-(a_+-K)}
\]
with the exact row-specific normalized overhead \(R_Q^{\rm row}\) fitting the remaining integer budget.  A proof through moments must satisfy the finite inequality
\[
        \log_2\Gamma_r+(a_+-K)\log_2|\B|\le r\Delta_Q
\]
with the actual allocated bit margin \(\Delta_Q\), not just an \(e^{o(n)}\) asymptotic loss.
\end{problem}

\begin{problem}[BC chart-decomposition target after moving-root]\label{prob:saturated-bc}
For each active MCA adjacent row, after quotient, boundary-Q, common-support, tangent, extension, degree-drop, and common-GCD branches have been paid by first match, verify that every remaining balanced-core split-pencil chart is one of the following two types:
\begin{enumerate}[label=\textup{(\alph*)}]
\item a projective one-parameter locator pencil covered by \cref{cor:bc-one-pencil}; or
\item a higher-dimensional coefficient family that is explicitly split into such pencils, or else assigned to a separate named residual cell with its own slope, not raw-support, bound.
\end{enumerate}
By \cref{prop:line-ray-saturation}, a raw support-locator census remains only an audit object,
\[
        \sum_{z\in E}\Cen(u+zv;a_+)
        =
        \sum_{(z,c)\in\LineRay_E(u,v;a_+)}\binom{s_{z,c}}{a_+},
\]
and can be exponentially larger than the slope numerator.  The enumerative obstruction for primitive one-parameter pencils is now removed by \cref{thm:bc-moving-root}; the finite row work is to certify that no residual high-dimensional BC family has escaped the pencil decomposition.
\end{problem}

\begin{remark}[finite arithmetic artifacts]\label{rem:finite-arithmetic-artifacts}
The exact integer tables in this note can be checked independently from the displayed formulas using integer arithmetic only: \(B^*=\lfloor p^s2^{-\lambda}\rfloor\), \(L(a)=\lceil\binom na p^{-(a-K)}\rceil\), the simple-pole value \(M(a)=\lceil L(a)(q-n)/(q-n+k(L(a)-1))\rceil\), the extension headroom in \cref{prop:extension-cell-target}, and the logarithmic Q-packing estimate in \cref{prop:proper-q-gap}, and the two-slope primitive one-pencil BC count in \cref{cor:bc-one-pencil}.  These checks are arithmetic audits; they do not replace the missing row-sharp Q theorem or the chart-decomposition verification in \cref{prob:row-sharp-q,prob:saturated-bc}.
\end{remark}

\section{Conditional Asymptotic Closure}
\label{sec:conditional-asymptotic-closure}

The previous sections isolate Q as the primitive max-fiber input and prove that SP is downstream from Q.  This section records the exact conditional implication from the additive-combinatorics atom \cref{prob:entropy-inverse-q} to the asymptotic RS--MCA frontier.  It is a theorem about the compiler once that atom is supplied; it is not a proof of the atom.

\begin{definition}[frontier sequence for MCA]\label{def:frontier-mca-sequence}
Let
\[
        C_n=\RS[\F_n,D_n,k_n],
        \qquad
        |D_n|=n,
        \qquad
        k_n/n\to\rho,
\]
be a sequence of smooth-domain Reed--Solomon rows with base fields \(\B_n\), and put \(K_n=k_n+1\) on the MCA route.  For an agreement threshold \(a_n\), set
\[
        m_n=a_n,
        \qquad
        w_n=a_n-K_n,
        \qquad
        \overline N_n=\binom n{m_n}|\B_n|^{-w_n}.
\]
The row is a frontier row if \(\overline N_n=\exp(o(n))\), and an extended frontier row if \(\overline N_n\ge\exp(-o(n))\).  A primitive residual leaf means the family remaining after quotient, planted, tangent, extension, Chebyshev, differential-locator, saturation, and split-pencil paid cells have been assigned by first match.
\end{definition}

\begin{lemma}[largest primitive fiber gives logarithmic excess]\label{lem:largest-fiber-to-excess-closure}
Let \(N^{\rm prim}(s)\), \(R^{\rm prim}(s)\), \(M_{\rm prim}\), and \(\Gamma_r^{\rm prim}\) be as in \cref{def:primitive-logmoment}.  Assume \(r\to\infty\) and \(w\log|\B|/r=o(n)\).  If \(M_{\rm prim}\ge\exp(\eta n)\) along a subsequence for some fixed \(\eta>0\), then, after replacing \(\eta\) by \(\eta/2\),
\[
        \Gamma_r^{\rm prim}\ge\exp((\eta/2)nr)
\]
along that subsequence.
\end{lemma}

\begin{proof}
Let \(s_0\) have \(R^{\rm prim}(s_0)=M_{\rm prim}\).  One term of the moment sum gives
\[
        \Gamma_r^{\rm prim}\ge |\B|^{-w}M_{\rm prim}^r
        \ge \exp(\eta nr-w\log|\B|).
\]
The hypothesis \(w\log|\B|/r=o(n)\) makes the exponent at least \((\eta/2)nr\) for all sufficiently large \(n\) in the subsequence.
\end{proof}

\begin{theorem}[primitive Q from the entropic inverse theorem]\label{thm:primitive-q-closure-module}
Assume that \cref{prob:entropy-inverse-q} holds uniformly for the primitive first-match residual leaves in an extended frontier sequence.  Then every primitive residual leaf satisfies
\[
        \max_s N^{\rm prim}(s)
        \le
        \exp(o(n))\binom n{m}|\B|^{-w}.
\]
Equivalently, \(M_{\rm prim}\le\exp(o(n))\).
\end{theorem}

\begin{proof}
If the conclusion failed, then for some fixed \(\eta>0\), \(M_{\rm prim}\ge\exp(\eta n)\) on an infinite subsequence.  Choose a logarithmic moment order \(r\to\infty\) with \(w\log|\B|/r=o(n)\).  By \cref{lem:largest-fiber-to-excess-closure}, that subsequence has positive-rate logarithmic excess.  Applying \cref{prob:entropy-inverse-q} gives either a paid algebraic cell or a positive-density rank defect among primitive Vandermonde columns.  The paid-cell alternative is excluded by the definition of the primitive residual leaf, and the rank-defect alternative contradicts \cref{prop:vandermonde-kills-low-rank}.  Hence no positive-rate failure of \(M_{\rm prim}\le\exp(o(n))\) remains.
\end{proof}

\begin{lemma}[subexponential add-back]\label{lem:subexponential-addback-closure}
Suppose that, in every first-match leaf, the number of paid quotient/profile cells is \(\exp(o(n))\), each paid cell contributes at most \(\exp(o(n))\overline N_n\) supports to any prefix syndrome, and the primitive residual satisfies \cref{thm:primitive-q-closure-module}.  Then the full residual prefix-boundary fiber satisfies
\[
        \max_sN(s)\le\exp(o(n))\overline N_n.
\]
\end{lemma}

\begin{proof}
For a fixed syndrome, decompose the fiber into first-match paid pieces and the primitive residual.  The primitive residual has size \(\exp(o(n))\overline N_n\) by \cref{thm:primitive-q-closure-module}.  There are \(\exp(o(n))\) paid pieces and each has the same scale.  Multiplying subexponential factors remains subexponential.
\end{proof}

\begin{proposition}[Q discharges SP at the frontier]\label{prop:q-sp-frontier-closure}
If \(\max_sN(s)\le\kappa_n\overline N_n\), then
\[
        \binom n{m}^{-1}\sum_sN(s)^2\le\kappa_n\overline N_n.
\]
In particular, if \(\kappa_n=\exp(o(n))\), the shift-pair ledger has the same asymptotic size as the Q max-fiber ledger.
\end{proposition}

\begin{proof}
Since \(\sum_sN(s)=\binom nm\),
\[
        \sum_sN(s)^2
        \le
        (\max_sN(s))\sum_sN(s)
        \le
        \kappa_n\overline N_n\binom nm .
\]
Dividing by \(\binom nm\) proves the claim.
\end{proof}

\begin{theorem}[asymptotic RS--MCA closure from the primitive inverse theorem]\label{thm:asymptotic-rs-mca-closure-combined}
Consider an extended frontier smooth-domain Reed--Solomon MCA sequence as in \cref{def:frontier-mca-sequence}.  Suppose the first-match ledger covers every MCA-bad finite slope by paid cells, residual Q, or residual SP; suppose the paid cells have total contribution \(\exp(o(n))\overline N_n\); suppose the quotient/profile add-back hypotheses of \cref{lem:subexponential-addback-closure} hold; and suppose \cref{prob:entropy-inverse-q} is true for every primitive residual leaf.  Then
\[
        B_{C_n}^{\rm MCA}(a_n)\le \exp(o(n))\overline N_n.
\]
Consequently, if the same hypotheses hold uniformly for agreements \(a_n=K_n+g_nn+o(n)\) on the safe side of the entropy--subfield crossing, then
\[
        \delta^*_{C_n}(\varepsilon^*)
        =
        1-\rho-g^*(\rho,\log_2|\B_n|)+o(1)
\]
whenever the identity-prefix lower construction supplies the matching unsafe side below the crossing.
\end{theorem}

\begin{proof}
The inverse theorem gives primitive Q by \cref{thm:primitive-q-closure-module}.  The add-back lemma upgrades primitive Q to the full residual Q ledger.  Then \cref{prop:q-sp-frontier-closure} bounds SP at the same asymptotic scale.  The paid part has scale \(\exp(o(n))\overline N_n\) by hypothesis, and first-match disjointness is exactly \cref{lem:first-match-ledger}.  Summing paid, Q, and SP cells gives the displayed numerator bound.

For the threshold identity, combine this upper bound with \cref{lem:entropy-bookkeeping}: above the entropy--subfield crossing, any \(o(n)\) agreement reserve whose entropy gain dominates the subexponential loss makes the safe numerator fit the target.  Below the crossing, \cref{thm:unsafe-envelope} supplies the identity-prefix unsafe numerator.  The remaining corridor has width \(o(n)\), which is the asserted \(o(1)\) radius error.
\end{proof}

\begin{remark}[why this is not a no-input proof]\label{rem:not-no-input-proof}
\Cref{thm:asymptotic-rs-mca-closure-combined} proves the compiler implication after the primitive entropy-scale inverse theorem has been supplied.  Removing that hypothesis requires proving \cref{prob:entropy-inverse-q}: a Renyi-scale trade-entropy bridge, an entropy BSG/PFR structural step or free-energy decay alternative, and the slice derivative from small-doubling trades back to a positive-density rank defect among individual Vandermonde columns.  Without that theorem, the adjacent finite rows and the asymptotic equality remain conditional safe-side statements.
\end{remark}

\section{Remaining Finite Upper-Ledger Target}

The proper proofs above remove the fake certificate-theorem closure.  They prove real replacement statements, but they do not settle the dense frontier.  A finite adjacent proof must improve the upper bound \(U(a_0+1)\) in \cref{thm:finite} until it fits inside the exact margins \(22.1969,22.0109,3.2589,3.0730\) bits.  After \cref{thm:q-implies-sp} and \cref{thm:bc-moving-root}, the independent active MCA targets are row-sharp multilevel Q and the finite BC chart-decomposition audit.  SP remains an exact audit ledger and follows once Q is strong enough; a primitive one-parameter BC pencil is now paid by the moving-root theorem.

\paragraph{Prefix-boundary ledger.}  The prefix-boundary row must bound the actual maximum fiber, not the raw null fiber.  \Cref{prop:mode-null-false} rules out the mode-at-null shortcut, and \cref{thm:moment-q} gives the exact moment inequality required for a finite proof:
\[
        \log_2\Gamma_r+w\log_2|\B|\le r\Delta .
\]
The new unconditional packing theorem \cref{thm:q-proper} is valid but, by \cref{prop:proper-q-gap}, is far too weak at the deployed adjacent rows.

\paragraph{Split-pencil interior ledger.}  The split-pencil proof must count saturated primitive line-rays, not raw support locators.  The exact saturation identity \cref{thm:saturation} explains the difference, and \cref{cor:raw-bc-fails} shows why the raw count can be exponentially larger than the MCA numerator.  The actual one-parameter split condition is now used in \cref{thm:bc-moving-root}: after common-GCD roots are paid, each moving root chooses at most one pencil parameter, so a primitive deployed BC pencil contributes at most two slopes.  Thus BC is no longer an enumerative conjecture for one-parameter locator pencils.  The remaining finite work is structural: the row ledger must show that every residual balanced-core chart is reduced to these pencils, or must isolate a genuine higher-dimensional residual cell.  The base-field floor in \cref{prop:base-field-floor} still warns that a challenge-field-only raw support normalization is false.

\paragraph{Primitive shift-pair ledger.}  The primitive pair ledger should remain in the paper as the exact \(L^2\) identity of \cref{prop:gamma2-ledger}.  It is no longer an independent active assumption in the MCA closure: \cref{thm:q-implies-sp} proves that a max-fiber Q theorem bounds the whole SP ledger.  The unconditional count in \cref{thm:sp-proper} is only an all-pairs audit ceiling.

\paragraph{Rung audit.}  The exact divisor-lattice audit at \(a_0+1\) should cover all quotient scales.  If a periodic rung exceeds the budget, the adjacent candidate fails for a paid reason and the edge moves.  If every rung is below budget with printed margin, the remaining proof target is genuinely primitive.

\section{Conclusion}

This version contains no finite Q, BC, or SP certificate theorem.  The unsafe side is exact and the broad safe anchors are unconditional in their stated ranges.  The former SP input is now a theorem downstream from Q; the raw BC support count is replaced by exact saturation and line-ray identities; and primitive one-parameter BC pencils are paid by the moving-root theorem.  On the asymptotic side, \cref{thm:logmoment-equivalence} proves that primitive Q is exactly a logarithmic collision-moment theorem under \(w\log|\B|/r=o(n)\), \cref{prob:entropy-inverse-q} isolates the standalone additive-combinatorics atom as an inverse theorem for Vandermonde slice sums, and \cref{thm:asymptotic-rs-mca-closure-combined} proves the conditional entropy--subfield threshold formula once that atom and the printed first-match ledger hypotheses are supplied.  \Cref{thm:fourier-flat-q} records a theorem-level sufficient condition for Q in Fourier-flat leaves, but it is not a deployed adjacent certificate unless the explicit Fourier error fits the row's finite margin.  The attempted Q closure in \cref{prop:q-exact-target,prop:q-moment-order-floor,prop:bc-not-q,prop:q-not-closed,prob:entropy-inverse-q} identifies the exact missing atom theorem and rules out the current shortcuts: packing is far too weak, fixed moments cannot fit the finite rows, standard polynomial-scale inverse Littlewood--Offord does not directly apply at exponential fiber scale, SP is downstream from Q rather than upstream, and the paid BC theorem applies to one-parameter pencils rather than the high-dimensional boundary-Q slice.  A closed adjacent theorem still requires a row-sharp Q atom bound together with the BC chart-decomposition audit required by \cref{prob:saturated-bc}.

\begin{thebibliography}{BCHKS25}

\bibitem[PP26]{ProximityPrize}
The Proximity Prize.
\newblock The Proximity Prize: \$1M Reed--Solomon Challenge.
\newblock Preliminary release, 2026.
\newblock \url{https://proximityprize.org/}.

\bibitem[ABF26]{ABF26}
G.~Arnon, D.~Boneh, and G.~Fenzi.
\newblock Remaining questions in list decoding and correlated agreement.
\newblock IACR Cryptology ePrint Archive, Report 2026/680, 2026.
\newblock \url{https://eprint.iacr.org/2026/680}.

\bibitem[BCHKS25]{BCHKS25}
E.~Ben-Sasson, D.~Carmon, U.~Hab{\"o}ck, S.~Kopparty, and S.~Saraf.
\newblock On proximity gaps for Reed--Solomon codes.
\newblock IACR Cryptology ePrint Archive, Report 2025/2055, 2025.
\newblock \url{https://eprint.iacr.org/2025/2055}.

\bibitem[BCIKS20]{BCIKS20}
E.~Ben-Sasson, D.~Carmon, Y.~Ishai, S.~Kopparty, and S.~Saraf.
\newblock Proximity gaps for Reed--Solomon codes.
\newblock In \emph{Proc. 61st IEEE Symposium on Foundations of Computer Science (FOCS)}, 2020.
\newblock Extended manuscript: IACR Cryptology ePrint Archive, Report 2020/654.

\bibitem[Cho26a]{Cho26a}
P.~Chojecki.
\newblock Capacity-edge obstructions to Reed--Solomon mutual correlated agreement over smooth multiplicative domains.
\newblock Manuscript, 2026.

\bibitem[Cho26b]{Cho26b}
P.~Chojecki.
\newblock Slack, quotient cores, and the entropy gap for smooth-domain Reed--Solomon codes.
\newblock Manuscript, merged corrected edition, 2026.

\bibitem[CAP25v13]{CAP25v13}
P.~Chojecki and contributors.
\newblock \emph{CAP25 v13 raw manuscript}.
\newblock Experimental manuscript in this repository, \texttt{experimental/cap25\_cap\_v13\_raw.tex}, 2026.

\bibitem[CS25]{CS25}
E.~Crites and A.~Stewart.
\newblock On Reed--Solomon proximity gaps conjectures.
\newblock IACR Cryptology ePrint Archive, Report 2025/2046, 2025.
\newblock \url{https://eprint.iacr.org/2025/2046}.

\bibitem[Lin44]{Linnik44}
Y.~V. Linnik.
\newblock On the least prime in an arithmetic progression. I. The basic theorem.
\newblock \emph{Rec. Math. [Mat. Sbornik] N.S.}, 15(57):139--178, 1944.

\bibitem[Wei48]{Weil48}
A.~Weil.
\newblock On some exponential sums.
\newblock \emph{Proceedings of the National Academy of Sciences of the United States of America}, 34(5):204--207, 1948.

\bibitem[LW10]{LiWanSieve}
J.~Li and D.~Wan.
\newblock A new sieve for distinct coordinate counting.
\newblock \emph{Science China Mathematics}, 53(9):2351--2362, 2010.

\bibitem[Tao05]{Tao05}
T.~Tao.
\newblock An uncertainty principle for cyclic groups of prime order.
\newblock \emph{Mathematical Research Letters}, 12(1):121--127, 2005.
\newblock Preprint: \url{https://arxiv.org/abs/math/0308286}.

\bibitem[TV09]{TaoVu09}
T.~Tao and V.~Vu.
\newblock A sharp inverse Littlewood--Offord theorem.
\newblock \emph{Random Structures \& Algorithms}, 37(4):525--539, 2010.
\newblock Preprint: \url{https://arxiv.org/abs/0902.2357}.

\bibitem[GR07]{GreenRuzsa07}
B.~Green and I.~Z. Ruzsa.
\newblock Freiman's theorem in an arbitrary abelian group.
\newblock \emph{Journal of the London Mathematical Society}, 75(1):163--175, 2007.
\newblock Preprint: \url{https://arxiv.org/abs/math/0505198}.

\bibitem[GGMT24]{GGMT24}
W.~T. Gowers, B.~Green, F.~Manners, and T.~Tao.
\newblock Marton's conjecture in abelian groups with bounded torsion.
\newblock Preprint, 2024.
\newblock \url{https://arxiv.org/abs/2404.02244}.

\end{thebibliography}

\end{document}
